\begingroup
\section{Node JC: honest joint endpoint comparison}
\label{module:JC}
\noindent\textit{Audited source: \nolinkurl{JC_joint_comparison.tex}.}
\subsection{Two honest scalars on one fixed fan}

Fix a positive normal fan $\alpha$, let $h_\alpha\mathbf1$ be its
area-$\pi$ tangential support vector, and work on the fixed affine
scale/translation section $\mathsf X_\alpha$ of the JP node.  Put
\begin{equation}\label{JC:eq:polygon-scalar}
 \Phi_\alpha(h)=\frac{|P_\alpha(h)|}{\pi}
                    \lambda_1(P_\alpha(h)).
\end{equation}
The closed JP--NS and disk-endpoint rows define
\begin{equation}\label{JC:eq:disk-data}
 B_\alpha=\Pi_{\rm ng}
 \{h_\alpha\sec(\theta-\theta_i)-1\},\qquad
 K_\alpha z=\Pi_{\rm ng}\mathcal I_\alpha z,
\end{equation}
and
\begin{equation}\label{JC:eq:metric-forcing}
 \Gjp_\alpha(z)=q_D(K_\alpha z),\qquad
 p_\alpha[z]=2\langle B_\alpha,K_\alpha z\rangle_D.
\end{equation}
Let $\mathcal R_\alpha$ be the Riesz map of $\Gjp_\alpha$ and set
\begin{equation}\label{JC:eq:disk-corrector}
 z_\alpha^D=-\frac12\mathcal R_\alpha^{-1}p_\alpha.
\end{equation}

\begin{proposition}[Exact disk completion]\label{JC:prop:disk-completion}
On $\mathsf X_\alpha$ one has
\begin{align}
 q_D(B_\alpha+K_\alpha z)
 &=P_D^{\rm ns}(\alpha)
   +\Gjp_\alpha(z-z_\alpha^D),\label{JC:eq:disk-completion}\\
 P_D^{\rm ns}(\alpha)
 &=q_D(B_\alpha)-\frac14
       \|p_\alpha\|_{(\Gjp_\alpha)^*}^2,\label{JC:eq:disk-schur}\\
 \Gjp_\alpha(z_\alpha^D)
 &=\frac14\|p_\alpha\|_{(\Gjp_\alpha)^*}^2.
 \label{JC:eq:corrector-size}
\end{align}
At the regular fan $\alpha_*$, cyclic symmetry gives
$p_{\alpha_*}=0$ and $z_{\alpha_*}^D=0$.
\end{proposition}

\begin{proof}
Expand the fixed quadratic form:
\[
 q_D(B_\alpha+K_\alpha z)
 =q_D(B_\alpha)+p_\alpha[z]+\Gjp_\alpha(z).
\]
The definition \eqref{JC:eq:disk-corrector} is exactly
$2\mathcal R_\alpha z_\alpha^D=-p_\alpha$.  Completing the Hilbert square
proves all three identities.  The regular assertion follows because the
regular source is cyclic while every vector in the quotient is orthogonal
to the cyclic mode.
\end{proof}

Define the fixed-fan joint residual
\begin{equation}\label{JC:eq:joint-residual}
\boxed{
 \mathscr R_\alpha(z):=
 \Phi_\alpha(h_\alpha\mathbf1+z)
 -\Phi_\alpha(h_\alpha\mathbf1)
 -\{q_D(B_\alpha+K_\alpha z)-q_D(B_\alpha)\}.}
\end{equation}
Both terms in \eqref{JC:eq:joint-residual} are genuine scalar functions of
$z$.  The disk term is the fixed Fourier--Bessel quadratic form; it is not
the eigenvalue of a moving disk domain.

\begin{proposition}[Exact integrated form]\label{JC:prop:integrated-form}
Let $\ell_\alpha=D_h\Phi_\alpha(h_\alpha\mathbf1)$ and
$e_\alpha=\ell_\alpha-p_\alpha$.  Whenever the fixed-normal support
segment is active,
\begin{equation}\label{JC:eq:integrated-form}
 \mathscr R_\alpha(z)=e_\alpha[z]
 +\int_0^1(1-t)
 \{D_h^2\Phi_\alpha(h_\alpha\mathbf1+tz)[z,z]
       -2\Gjp_\alpha(z)\}\,dt.
\end{equation}
\end{proposition}

\begin{proof}
Apply the fundamental theorem of calculus twice to the polygonal scalar.
For the disk scalar use
$q_D(B+Kz)-q_D(B)=p[z]+\Gjp(z)$.  Subtraction gives
\eqref{JC:eq:integrated-form} exactly.
\end{proof}

\subsection{Exact Schur reduction and support-time cancellation}

Let $g_\alpha(\cdot,\cdot)$ be the bilinear form associated with
$\Gjp_\alpha$, put $d_\alpha=z_\alpha^D$, and abbreviate
\begin{equation}\label{JC:eq:Hessian-family}
 H_{\alpha,y}=D_h^2\Phi_\alpha(h_\alpha\mathbf1+y).
\end{equation}
Equation~\eqref{JC:eq:disk-corrector} is equivalently
\begin{equation}\label{JC:eq:p-d-relation}
 p_\alpha[\xi]=-2g_\alpha(d_\alpha,\xi).
\end{equation}
Define the completed covector
\begin{equation}\label{JC:eq:completed-covector}
 b_\alpha[\xi]
 :=\ell_\alpha[\xi]+g_\alpha(d_\alpha,\xi)
 =\ell_\alpha[\xi]-\frac12p_\alpha[\xi]
\end{equation}
and the pathwise Schur action
\begin{equation}\label{JC:eq:pathwise-action}
 \mathscr A_\alpha(z):=
 b_\alpha[z]
 +\int_0^1(1-t)
 \{H_{\alpha,tz}[z,z]-g_\alpha(z,z)\}\,dt
 +\frac12g_\alpha(d_\alpha,d_\alpha)
 +\frac\kappa2\Jfour(\alpha).
\end{equation}

\begin{proposition}[JC--ALG: exact completed-action equivalence]
\label{JC:prop:JC-ALG}
For every active fixed-fan segment,
\begin{equation}\label{JC:eq:JC-equivalence}
 \boxed{
 \mathscr A_\alpha(z)
 =\mathscr R_\alpha(z)
  +\frac12\Gjp_\alpha(z-z_\alpha^D)
  +\frac\kappa2\Jfour(\alpha).}
\end{equation}
Consequently the former JP--JC inequality is equivalent, without a loss
of constants, to $\mathscr A_\alpha(z)\ge0$.
\end{proposition}

\begin{proof}
Insert the integrated identity~\eqref{JC:eq:integrated-form} and expand
\[
 \frac12g_\alpha(z-d_\alpha,z-d_\alpha)
 =\frac12g_\alpha(z,z)-g_\alpha(z,d_\alpha)
   +\frac12g_\alpha(d_\alpha,d_\alpha).
\]
By~\eqref{JC:eq:p-d-relation},
$-g_\alpha(z,d_\alpha)=p_\alpha[z]/2$.  Hence the linear row is
$(\ell_\alpha-p_\alpha)[z]+p_\alpha[z]/2=b_\alpha[z]$.
Moreover
\[
 \int_0^1(1-t)(-2g_\alpha(z,z))\,dt
 +\frac12g_\alpha(z,z)=-\frac12g_\alpha(z,z),
\]
which is exactly the integral of $-g_\alpha(z,z)$ with weight $1-t$.
The constant rows agree with~\eqref{JC:eq:pathwise-action}, proving
\eqref{JC:eq:JC-equivalence}.
\end{proof}

The same identity has the following equivalent form, which displays the
positive disk corrector square before any estimate:
\begin{equation}\label{JC:eq:joint-Peano-block}
 \boxed{
 \mathscr A_\alpha(z)
 =\frac12\Gjp_\alpha(z-d_\alpha)
  +e_\alpha[z]
  +\int_0^1(1-t)
    \{H_{\alpha,tz}[z,z]-2\Gjp_\alpha(z)\}\,dt
  +\frac\kappa2\Jfour(\alpha).}
\end{equation}
Indeed, $b_\alpha=e_\alpha-g_\alpha(d_\alpha,\cdot)$ and
$\int_0^1(1-t)g_\alpha(z,z)\,dt=\Gjp_\alpha(z)/2$.
Thus the period-three component of $e_\alpha$ is not to be bounded as an
isolated covector: it must be paid jointly by the jump surplus in
$H_{\alpha,tz}-2\Gjp_\alpha$ and the fan row.

\begin{lemma}[JC--TL: support-time logarithm cancellation]
\label{JC:lem:support-time-log}
Let $\beta>0$, $0<r\le1$, and
\begin{equation}\label{JC:eq:I-beta}
 I_\beta(r):=\int_0^1(1-t)
 \left\{1-\bigl((1-t)+tr\bigr)^{2\beta}\right\}\,dt.
\end{equation}
Then
\begin{align}
 0\le I_\beta(r)
 &\le\frac\beta3\log\frac1r,\label{JC:eq:log-bound}\\
 I_\beta(r)
 &\le\frac{\beta}{2(1+\beta)}\le\frac\beta2.\label{JC:eq:absolute-bound}
\end{align}
In particular,
\begin{equation}\label{JC:eq:min-bound}
 \boxed{I_\beta(r)\le\frac\beta3
 \min\left\{\log\frac1r,\frac32\right\}.}
\end{equation}
All constants are independent of $r$.
\end{lemma}

\begin{proof}
For $0<x\le1$, $1-x^{2\beta}\le2\beta(-\log x)$.
Weighted arithmetic--geometric mean gives
$(1-t)+tr\ge r^t$ and hence
\[
 -\log((1-t)+tr)\le t\log(1/r).
\]
Since $\int_0^1t(1-t)\,dt=1/6$, these two inequalities prove
\eqref{JC:eq:log-bound}.  Also $(1-t)+tr\ge1-t$, so monotonicity of the
power gives
\[
 I_\beta(r)\le I_\beta(0)
 =\frac12-\frac1{2+2\beta}
 =\frac\beta{2(1+\beta)}.
\]
This proves~\eqref{JC:eq:absolute-bound}; taking the smaller bound gives
\eqref{JC:eq:min-bound}.
\end{proof}

\begin{lemma}[JC--TL$^{\min}$: switches of the shortest incident side]
\label{JC:lem:support-time-min}
Let $L_-$ and $L_+$ be positive affine functions on $[0,1]$, put
$m(t):=\min\{L_-(t),L_+(t)\}$, and set $r:=m(1)/m(0)$.  Then, for every
$\beta>0$,
\begin{equation}\label{JC:eq:min-affine-payment}
 \int_0^1(1-t)\left\{1-
 \left(\frac{m(t)}{m(0)}\right)^{2\beta}\right\}\,dt
 \le
 \begin{cases}
 I_\beta(r),&0<r\le1,\\
 0,&r\ge1.
 \end{cases}
\end{equation}
Consequently every \emph{upper} bound in
Lemma~\ref{JC:lem:support-time-log} remains valid for the positive part of the
loss when a single affine length is replaced by the shortest of the two
incident physical lengths.  No switch-point boundary term is present.
\end{lemma}

\begin{proof}
The pointwise minimum of affine functions is concave.  If $r\le1$,
concavity gives
\[
 \frac{m(t)}{m(0)}\ge(1-t)+tr.
\]
The function $x\mapsto1-x^{2\beta}$ is decreasing, so integration with
the nonnegative weight $1-t$ proves the first line of
\eqref{JC:eq:min-affine-payment}.  If $r\ge1$, the chord joining the endpoint
values lies above $m(0)$, and concavity gives $m(t)\ge m(0)$ throughout;
the integrand is therefore nonpositive.  The argument differentiates
neither $m$ nor $m^{2\beta}$, hence a change of the minimizing affine side
creates no distributional boundary contribution.
\end{proof}

\begin{remark}[Exact terminal normalization]
For an actual physical side length with the affine form
$L(t)=L_0\{(1-t)+tr\}$, the squared corner-flux power uses
\begin{equation}\label{JC:eq:physical-corner-exponent}
 \beta_i^{\rm ph}=\frac{\alpha_i}{\pi-\alpha_i},
\end{equation}
not $\alpha_i/\pi$.  Put
$\Delta_i=a_i^--a_i^+=2\sin(\alpha_i/2)\tau_i$.  The integrated terminal
amplitude loss obeys
\begin{equation}\label{JC:eq:terminal-log-payment}
 \alpha_i\tau_i^2 I_{\beta_i^{\rm ph}}(r)
 \le |\Delta_i|^2\min\left\{
 \frac{\alpha_i^2}{12(\pi-\alpha_i)\sin^2(\alpha_i/2)}
       \log\frac1r,
 \frac{\alpha_i^2}{8\pi\sin^2(\alpha_i/2)}\right\}.
\end{equation}
Thus the logarithmic coefficient is $1/(3\pi)+O(\alpha_i)$, while the
exact absolute cap is
\begin{equation}\label{JC:eq:terminal-absolute-payment}
 \alpha_i\tau_i^2 I_{\beta_i^{\rm ph}}(r)
 \le c_T(S)|\Delta_i|^2,
 \qquad
 c_T(S):=\sup_{0<\alpha\le S}
 \frac{\alpha^2}{8\pi\sin^2(\alpha/2)},
\end{equation}
and $c_T(S)\to1/(2\pi)$ as $S\downarrow0$.  This is the sharp
support-time cap for a terminal length tending to zero.  A growing
terminal length has the favorable sign and is separated before the
absolute estimate.
The exponent $\alpha_i/\pi$ belongs instead to a conformal-prevertex
power.  No retained node proves that a conformal prevertex interval is
affine in support time, so that substitution is not allowed.  If the
terminal cutoff is the smaller of two incident physical lengths,
Lemma~\ref{JC:lem:support-time-min} applies globally: concavity of the
minimum dominates the affine endpoint chord, so neither subdivision nor a
switch correction is needed.  (Restarting the Taylor weight after a
subdivision would in fact change the ledger and is not the argument used
here.)
The right side is the descendant-annulus/terminal-cap jump charge on the
same physical scale.  Thus taking a side-ratio supremum before the
support-time integral destroys a real cancellation and is inadmissible.
\end{remark}

Formula \eqref{JC:eq:integrated-form} is the reason the two old obligations
must not be estimated independently.  In the three-periodic obstruction
$e_\alpha$ has a non-vanishing normalized leading term; the surplus in the
integrated Hessian is available in \eqref{JC:eq:integrated-form} to complete
the corresponding Schur square.

\begin{proposition}[JC--RS: exact quadratic root certificate]
\label{JC:prop:JC-RS}
Put
\begin{equation}\label{JC:eq:root-form}
 Q_\alpha[\xi,\xi]
 :=H_{\alpha,0}[\xi,\xi]-g_\alpha(\xi,\xi).
\end{equation}
Define the frozen quadratic root model
\begin{equation}\label{JC:eq:quadratic-root-model}
 \mathscr A_\alpha^{(2)}(\xi)
 :=\frac12g_\alpha(d_\alpha,d_\alpha)
   +\frac\kappa2\Jfour(\alpha)
   +b_\alpha[\xi]+\frac12Q_\alpha[\xi,\xi].
\end{equation}
Then $\mathscr A_\alpha^{(2)}(\xi)\ge0$ for every $\xi$ if and only if
\begin{equation}\label{JC:eq:root-Schur-necessary}
 \boxed{
 Q_\alpha\ge0,
 \qquad
 \|b_\alpha\|_{Q_\alpha^*}^2
 \le g_\alpha(d_\alpha,d_\alpha)
       +\kappa\Jfour(\alpha).}
\end{equation}
For a semidefinite $Q_\alpha$, the second statement includes that
$b_\alpha$ annihilates its kernel and uses the quotient dual norm.
\end{proposition}

\begin{proof}
This is Hilbert completion of the quadratic form
\eqref{JC:eq:quadratic-root-model}.  If it is nonnegative on the whole vector
space, its quadratic part is semidefinite and its linear part annihilates
the kernel.  Its infimum is
\[
 \frac12g_\alpha(d_\alpha,d_\alpha)
 +\frac\kappa2\Jfour(\alpha)
 -\frac12\|b_\alpha\|_{Q_\alpha^*}^2,
\]
which is nonnegative exactly under~\eqref{JC:eq:root-Schur-necessary}.  The
reverse implication follows from the same completed square.
\end{proof}

Proposition~\ref{JC:prop:JC-RS} is an exact certificate for the frozen
quadratic root model, not a necessary consequence of local nonnegativity
of the full action: its positive constant term prevents that inference.
To use it for JC--PL one must additionally compare the full integrated
Hessian to this frozen model on the whole admissible amplitude range.  The
closed JP--DE estimates control $p_\alpha$ and hence $d_\alpha$, but do not
prove either that comparison or the displayed sharp dual bound for
$b_\alpha=\ell_\alpha-p_\alpha/2$.  The period-three computation passes
the quadratic certificate with a sizeable numerical margin; a uniform
no-ratio proof over all profiles is still required.

\begin{proposition}[JC--RC: exact root/path sufficient reduction]
\label{JC:prop:JC-RC}
Suppose a positive semidefinite form $L_\alpha$ satisfies, throughout the
radial/metric path chart,
\begin{equation}\label{JC:eq:path-loss-form}
 \int_0^1(1-t)
 \{H_{\alpha,tz}-H_{\alpha,0}\}[z,z]\,dt
 \ge-L_\alpha[z,z].
\end{equation}
Define
\begin{equation}\label{JC:eq:root-after-path-loss}
 \widehat Q_\alpha:=H_{\alpha,0}-g_\alpha-2L_\alpha.
\end{equation}
If
\begin{equation}\label{JC:eq:root-path-Schur}
 \widehat Q_\alpha\ge0,
 \qquad
 \|b_\alpha\|_{\widehat Q_\alpha^*}^2
 \le g_\alpha(d_\alpha,d_\alpha)+\kappa\Jfour(\alpha),
\end{equation}
then JC--PL holds.
\end{proposition}

\begin{proof}
Split the Hessian integral in~\eqref{JC:eq:pathwise-action} at
$H_{\alpha,0}$.  The Bregman weight has mass $1/2$, so
\eqref{JC:eq:path-loss-form} gives
\[
 \mathscr A_\alpha(z)
 \ge\frac12\{g_\alpha(d_\alpha,d_\alpha)
                 +\kappa\Jfour(\alpha)\}
     +b_\alpha[z]
     +\frac12\widehat Q_\alpha[z,z].
\]
The right side is nonnegative by the same Hilbert completion as
Proposition~\ref{JC:prop:JC-RS}.
\end{proof}

This reduction records the required slack correctly.  Terminal and
nonterminal estimates must first be assembled into a genuine form
$L_\alpha$; one may not subtract their scalar upper bounds twice or
replace~\eqref{JC:eq:path-loss-form} by a pointwise support-time supremum.

\begin{remark}[Why the tempting exact Gram shortcut does not close the root]
Let $\mathcal H_D$ be the disk boundary Hilbert space, let $K$ be the
fan-normal sine trace, $T$ the genuine material/kinematic trace at the
tangential endpoint, and $B$ the fan source.  Put
\[
 G=K^*K,\qquad p=2K^*B,\qquad
 d=-G^{-1}p/2,\qquad c=-Kd=P_{K\mathsf X_\alpha}B,
 \qquad E=T-K.
\]
For the pure quadratic disk principal block, where
$\ell=2T^*B$ and $H_0=2T^*T$, direct expansion gives, with $V=2T-K$,
\begin{align}
 b[z]&=\langle c,Vz\rangle_D
       +2\langle B-c,Ez\rangle_D,
 \label{JC:eq:Gram-linear-defect}\\
 H_0-G&=V^*V-2E^*E,
 \label{JC:eq:Gram-quadratic-defect}\\
 \frac12G(d)+b[z]+\frac12(H_0-G)[z,z]
 &=\frac12\|c+Kz+2Ez\|_D^2-
   \|Ez\|_D^2+2\langle B-c,Ez\rangle_D.
 \label{JC:eq:Gram-failed-square}
\end{align}
Thus choosing $U=c$ and $V=2T-K$ leaves the negative form $-2E^*E$
and an orthogonal-source cross term; they do not vanish merely because
$c$ is the projection of $B$ onto $K\mathsf X_\alpha$.

For the genuine polygonal block write
$\Gamma=H_{\alpha,0}-2T^*T$ and let $r$ be the difference between
$b_\alpha$ and the first pairing in~\eqref{JC:eq:Gram-linear-defect}.  A
valid Gram repair must, at minimum, pay the block
\[
 \begin{pmatrix}
  \kappa\Jfour(\alpha) & r^*\\ r & \Gamma-2E^*E
 \end{pmatrix}
 \ge0
\]
after adding any explicitly retained fan square and path-loss slack.
No closed TR or JP--DE row proves this block.  Equations
\eqref{JC:eq:Gram-linear-defect}--\eqref{JC:eq:Gram-failed-square} explain why
the numerically favorable root ratio cannot be promoted by a bare
Cauchy--Schwarz citation.
\end{remark}

\begin{remark}[Audited physical-trace branch: exact algebra, failed chart]
There is an algebraically cleaner inactive branch.  Let
$\mathbf q_D=2j_{0,1}^2q_D^0$ be the physical disk half-Hessian, put
\[
 \widetilde K_\alpha z=\Pi_{\rm ng}\widetilde T_\alpha z,
 \qquad
 \widetilde G_\alpha(z)=\mathbf q_D(\widetilde K_\alpha z),
 \qquad
 \widetilde p_\alpha[z]
 =2\langle B_\alpha,\widetilde K_\alpha z\rangle_D,
\]
and let
$\widetilde d_\alpha=-\tfrac12\widetilde{\mathcal R}_\alpha^{-1}
\widetilde p_\alpha$.  The normalized TR endpoint gives the exact square
\[
 \mathbf q_D(B_\alpha+\widetilde K_\alpha z)
 =\widetilde P_D(\alpha)
  +\widetilde G_\alpha(z-\widetilde d_\alpha).
\]
The corresponding physical residual and action satisfy the same exact
Schur identity as Proposition~\ref{JC:prop:JC-ALG}; in its pure disk block
the root is exactly
\[
 \frac12\widetilde G_\alpha(z-\widetilde d_\alpha)
 +\frac\kappa2\Jfour(\alpha)\ge0.
\]
Thus the $E=T-K$ defect in~\eqref{JC:eq:Gram-quadratic-defect} disappears
from the disk block.

This does not close the polygonal action.  The large component of
$\widetilde T_\alpha$ produced by a short angle is a tangential slide of a
straight side viewed in the disk radial frame; it can be nonzero even
when the true side-normal speed is zero.  It is cancelled at the same
order by the nonlinear reparametrization remainder $\mathcal N$:
\[
 D_z^2\widetilde{\mathscr R}_\alpha(0)[\xi,\xi]
 =H_{\alpha,0}[\xi,\xi]-2\widetilde G_\alpha(\xi).
\]
By the exact value split in TR, the difference on the right is precisely
the full second $z$-derivative of $\mathcal N(W_{\alpha,z})$ at $z=0$;
this formulation includes all chain-rule terms at the vector-field base
point $W_{\alpha,0}$.
The cancellation cannot be estimated from the radial image alone.  For
example, if
$X(\theta)=e_r(\theta+\delta\sin2\theta)$ and $W=X-e_r$, the boundary
image is the unit circle, so $\mathcal L(W)=\mathcal L(0)$, whereas
\[
 f=W\cdot e_r=\cos(\delta\sin2\theta)-1,
 \qquad
 \mathcal N(W)=-\mathbf q_D(\Pi_{\rm ng}f).
\]
Hence no image-small argument can yield
$\mathcal N(W)\ge-\vartheta\mathbf q_D(\Pi_{\rm ng}f)$ with a universal
$\vartheta<1$.  The radial-not-affine example in JP further shows that
the full material $W^{1,\infty}$ norm can diverge on the radial chart.
Consequently the old physical CF+CM$\to$NR support-tail argument cannot
be reactivated.  The physical branch replaces JC--PL by an equivalent
open full-radial-chart action; it is recorded as an inactive alternative,
not as a second proof of J0.
\end{remark}

\begin{remark}[Audited Minkowski shortcut: two concavities do not order the ratio]
For a fixed positive normal fan, support interpolation is Minkowski
interpolation.  Consequently
\[
 U(h)=|P(h)|^{1/2},
 \qquad V(h)=\lambda_1(P(h))^{-1/2}
\]
are positive, one-homogeneous concave functions, while
$\Phi=(U/V)^2$.  At a stationary point,
\[
 \frac{D^2\Phi}{\Phi}
 =2\left(\frac{D^2U}{U}-\frac{D^2V}{V}\right)
\]
on a nonradial direction.  Both displayed curvatures are nonpositive,
but their difference has no determined sign.

This is a genuine logical obstruction, not only a missing inequality.
On the Lorentz cone $\{(x,y):x>|y|\}$ put, with $c=11/20$,
\[
 U(x,y)=\sqrt{x^2-y^2},
 \qquad V(x,y)=x-c\frac{y^2}{x}.
\]
Both functions are positive, one-homogeneous and concave, since
\[
 D^2U=-\frac1{(x^2-y^2)^{3/2}}
 \begin{pmatrix}y^2&-xy\\-xy&x^2\end{pmatrix},
 \quad
 D^2V=-\frac{2c}{x^3}
 \begin{pmatrix}y^2&-xy\\-xy&x^2\end{pmatrix}.
\]
Moreover $U^2=x^2-y^2$ is exactly quadratic.  At $(1,0)$ the ratio
$(U/V)^2$ is stationary and has second derivative $1/5>0$ in the
$y$ direction.  Nevertheless the same-area point $(5/3,4/3)$ satisfies
\[
 U(5/3,4/3)=1,
 \qquad
 \left(\frac UV\right)^2(5/3,4/3)=\frac{625}{729}<1.
\]
Thus area and eigenvalue Brunn--Minkowski concavity, even with exact
quadratic area and strict local support minimality, do not give a global
support comparison.  The missing relative curvature comparison is again
part of JC--PL.
\end{remark}

\begin{remark}[Audited stationary-endpoint bypass: why it is insufficient]
The independent stationary O--H endpoint theorem cannot by itself replace
JC--PL.  Indeed, put the regular polygon and a hypothetical active global
minimizer in one fixed-area joint chart at $0$ and $x$.  Since both are
stationary, the exact identity is
\begin{equation}\label{JC:eq:stationary-gradient-chord}
 0=\langle\nabla\Phi_N(x)-\nabla\Phi_N(0),x\rangle
 =\int_0^1D^2\Phi_N(tx)[x,x]\,dt.
\end{equation}
The interior points $tx$ are not stationary, so a Hessian theorem whose
hypothesis is stationarity does not control this integral.

Endpoint coercivity alone is logically insufficient even at the relevant
super-near scale.  For $a>0$ and $0<c<1/3$ consider
\[
 F_{a,c}(t,y)=a t^2(a-t)^2
 -ca^5\{3(t/a)^2-2(t/a)^3\}+|y|^2,
 \qquad E_a(u,v)=a^3u^2+|v|^2.
\]
The points $(0,0)$ and $(a,0)$ are stationary, while
$F_{a,c}(a,0)=-ca^5<F_{a,c}(0,0)=0$.  Both endpoint Hessians are strictly
positive: their $t$-eigenvalues are respectively
\[
 (2-6c)a^3,\qquad (2+6c)a^3,
\]
and the $y$-block is $2I$.  Moreover, with $u=t/a$,
\[
 \frac{F_{a,c}(t,0)}{a^5}+c
 =(1-u)^2\{(u+c)^2+c-c^2\}\ge0,
\]
so $(a,0)$ is the global minimum.  Nevertheless
$D_{tt}^2F_{a,c}(a/2,0)=-a^3$ and, for $x=(a,0)$,
\[
 \int_0^1D^2F_{a,c}(sx)[x,x]\,ds=0,
 \qquad E_a(x)=a^5.
\]
Thus positive endpoint spectra, including uniform spectra in the naturally
scaled norm, do not prove the path monotonicity needed in
\eqref{JC:eq:stationary-gradient-chord}.

The minimal stationary-route replacement is a super-near joint statement,
not an endpoint citation.  After gauges, it would have to prove uniformly
\begin{equation}\label{JC:eq:stationary-super-near-cut}
 \int_0^1D^2\Phi_N(tx)[x,x]\,dt
 \ge\mu\left\{6a^2\sum_i d_i^2
 +\Gjp_{\alpha_*}(z-Z_Nd)\right\}
\end{equation}
for nonzero $x=(d,z)$ with
$\Jfour(\alpha_*+d)+\Gjp_{\alpha_*}(z)\lesssim a^5$.
At the regular polygon this requires the full cyclic fan/support Schur
symbol, including its nonzero period-three mixed block; it must then be
propagated diagonally along the super-near path.  No retained node proves
either the regular Schur complement or this relative path stability.
Consequently the stationary route only narrows JC--PL; it does not close
the main theorem.
\end{remark}

\subsection{The unique open pathwise payment}

The closed disk endpoint and value rows provide universal
$S_D,\kappa,C_p>0$ and a modulus $\varepsilon_V(S)\to0$ such that, for
$S\le S_D$,
\begin{align}
 P_D^{\rm ns}(\alpha)-P_D^{\rm ns}(\alpha_*)
 &\ge2\kappa\Jfour(\alpha),\label{JC:eq:D0}\\
 \|p_\alpha\|_{(\Gjp_\alpha)^*}
 &\le C_p\Jfour(\alpha)^{1/2},\label{JC:eq:P}\\
 \left|\Phi_\alpha(h_\alpha\mathbf1)
 -\Phi_{\alpha_*}(h_*\mathbf1)
 -\{q_D(B_\alpha)-q_D(B_{\alpha_*})\}\right|
 &\le\varepsilon_V(S)\Jfour(\alpha).
 \label{JC:eq:V}
\end{align}

\begin{proposition}[Regular disk-joint Schur symbol]
\label{JC:prop:regular-disk-joint}
Let $\alpha_*=a\mathbf1$ with $2a<S_D$.  Choose the canonical analytic
scale/translation trivialization constructed in
Proposition~\ref{JC:prop:super-near-bridge-active},
\[
 J_\alpha:\mathsf X_{\alpha_*}\longrightarrow\mathsf X_\alpha,
 \qquad J_{\alpha_*}=I,
\]
and put
\[
 F_D^\sharp(\alpha,z):=
 q_D(B_\alpha+K_\alpha J_\alpha z),\qquad
 \bar z_\alpha^D:=J_\alpha^{-1}z_\alpha^D,\qquad
 Z_N\delta:=
 D_\alpha\bar z_\alpha^D\big|_{\alpha_*}[\delta].
\]
For every $\sum_i\delta_i=0$ and every regular-quotient support vector
$\zeta$,
\begin{align}
 \mathcal D_N(\delta,\zeta)
 &:=
 \frac12D^2F_D^\sharp(\alpha_*,0)
 [(\delta,\zeta),(\delta,\zeta)]
 \label{JC:eq:regular-disk-joint-identity}\\
 &=
 \frac12D^2P_D^{\rm ns}(\alpha_*)[\delta,\delta]
 +\Gjp_{\alpha_*}(\zeta-Z_N\delta)
 \notag\\
 &\ge
 12\kappa a^2\sum_i\delta_i^2
 +\Gjp_{\alpha_*}(\zeta-Z_N\delta).
 \label{JC:eq:regular-disk-joint-gap}
\end{align}
Thus every cyclic Fourier block of the regular
\emph{normal-sine disk-joint} symbol has a uniform completed Schur gap.
This statement concerns the fixed disk scalar; it is not a lower bound
for the polygonal joint Hessian.
\end{proposition}

\begin{proof}
The endpoint data and $J_\alpha$ are analytic in the fixed quotient
chart.  If
$\bar\Gjp_\alpha(w):=\Gjp_\alpha(J_\alpha w)$, exact disk completion gives
\[
 F_D^\sharp(\alpha,z)
 =P_D^{\rm ns}(\alpha)
  +\bar\Gjp_\alpha(z-\bar z_\alpha^D),
\qquad \bar z_{\alpha_*}^D=0.
\]
Along $\alpha(t)=\alpha_*+t\delta$, $z(t)=t\zeta$, one has
$\bar z_{\alpha(t)}^D=tZ_N\delta+O(t^2)$, while variation of
$\bar\Gjp_{\alpha(t)}$ contributes only from order $t^3$ onward.  Taking
the coefficient of $t^2$ proves
\eqref{JC:eq:regular-disk-joint-identity}.  Moreover
\[
 \Jfour(\alpha_*+t\delta)
 =6a^2t^2\sum_i\delta_i^2
  +4at^3\sum_i\delta_i^3+t^4\sum_i\delta_i^4.
\]
Apply the closed estimate
$P_D^{\rm ns}(\alpha)-P_D^{\rm ns}(\alpha_*)
\ge2\kappa\Jfour(\alpha)$.  For fixed $\delta$, both signs of sufficiently
small $t$ keep $\alpha_*+t\delta$ in the positive-fan region
$S\le S_D$.  Dividing the two signed inequalities first by $t$ shows
$DP_D^{\rm ns}(\alpha_*)[\delta]=0$; subtracting that linear term and
then dividing by $t^2$ gives
$\frac12D^2P_D^{\rm ns}(\alpha_*)[\delta,\delta]
\ge12\kappa a^2\sum_i\delta_i^2$, proving
\eqref{JC:eq:regular-disk-joint-gap}.
\end{proof}

\begin{contract}[JC--AR: regular completed-action Schur block]
\label{JC:con:JC-AR}
In the trivialization of Proposition~\ref{JC:prop:regular-disk-joint}, put
\[
 F_P^\sharp(\alpha,z)
 :=\Phi_\alpha(h_\alpha\mathbf1+J_\alpha z),
 \qquad
 \mathcal E_N(d,\zeta)
 :=\frac12D^2(F_P^\sharp-F_D^\sharp)(\alpha_*,0)
 [(d,\zeta),(d,\zeta)].
\]
Write \(\bar z_\alpha^D=J_\alpha^{-1}z_\alpha^D\), so that
\(D\bar z_{\alpha_*}^D[d]=Z_Nd\), and use the exact completed coordinates
\[
 \mathbf A_N(d,w):=
 \mathscr A_{\alpha_*+d}
 \bigl(J_{\alpha_*+d}(\bar z_{\alpha_*+d}^D+w)\bigr).
\]
Then \(\mathbf A_N(0,0)=0\), \(D\mathbf A_N(0,0)=0\), and the exact
regular quadratic root is
\begin{equation}\label{JC:eq:regular-action-root}
 \boxed{
 \mathcal Q_N(d,w):=\frac12D^2\mathbf A_N(0,0)[(d,w),(d,w)]
 =\mathcal E_N(d,w+Z_Nd)-\mathcal E_N(d,0)
 +\frac12\Gjp_{\alpha_*}(w)
 +3\kappa a^2\sum_i d_i^2.}
\end{equation}
The subtraction \(\mathcal E_N(d,0)\) and the coefficients
\(1/2\) and \(3\kappa\) are forced respectively by the base-value
subtraction in \(\mathscr R_\alpha\), the completed disk square, and
\(\frac\kappa2\Jfour\).  They may not be replaced by the larger margins
of the full disk-joint Hessian.

Let
\[
 \mathcal P_N(d,\zeta):=
 \frac12D^2F_P^\sharp(\alpha_*,0)
 [(d,\zeta),(d,\zeta)]
\]
be the genuine polygonal joint half-Hessian.  Exact disk completion gives
the following sharper algebraic reduction:
\begin{equation}\label{JC:eq:regular-action-polygon-reduction}
 \boxed{
 \mathcal Q_N(d,w)=
 \mathcal P_N(d,w+Z_Nd)-\mathcal P_N(d,0)
 +\Gjp_{\alpha_*}(Z_Nd)-\frac12\Gjp_{\alpha_*}(w)
 +3\kappa a^2\sum_i d_i^2.}
\end{equation}
Indeed Proposition~\ref{JC:prop:regular-disk-joint} yields
\[
 \frac12D^2F_D^\sharp(\alpha_*,0)[(d,\zeta)]^2
 =\frac12D^2P_D^{\rm ns}(\alpha_*)[d,d]
  +\Gjp_{\alpha_*}(\zeta-Z_Nd).
\]
Subtracting its values at $\zeta=w+Z_Nd$ and $\zeta=0$ gives
$\Gjp(w)-\Gjp(Z_Nd)$, which proves
\eqref{JC:eq:regular-action-polygon-reduction}.  In particular the pure
polygonal fan block cancels exactly.  JC--AR therefore asks only for the
polygonal support block and its fan--support coupling after the corrector
is inserted; estimating the entire polygonal joint Hessian is unnecessary.

For $d=0$, equation~\eqref{JC:eq:regular-action-polygon-reduction} becomes
\begin{equation}\label{JC:eq:regular-action-pure-support-identity}
 \mathcal Q_N(0,w)=
 \frac12D_h^2\Phi_{\alpha_*}(h_*\mathbf1)[w,w]
 -\frac12\Gjp_{\alpha_*}(w).
\end{equation}
Consequently the archived regular specialization of JP--S0,
\[
 D_h^2\Phi_{\alpha_*}(h_*\mathbf1)[w,w]
 \ge(2-\omega_0(a))\Gjp_{\alpha_*}(w),
 \qquad \omega_0(a)\to0,
\]
would close the pure-support row of JC--AR with any fixed
$\mu_s<1/2$ for all sufficiently large $N$.  The remaining part is then
the modewise Schur payment of the fan--support coupling in
\eqref{JC:eq:regular-action-polygon-reduction}.

The open regular action target is to prove constants
\begin{equation}\label{JC:eq:regular-action-constants}
 \mu_f>0,\qquad \mu_s>0,
\end{equation}
independent of $N$, such that
\begin{equation}\label{JC:eq:regular-action-Schur}
 \boxed{
 \mathcal Q_N(d,w)
 \ge \mu_fa^2\sum_i d_i^2+
      \mu_s\Gjp_{\alpha_*}(w).}
\end{equation}
This is not a consequence of Proposition~\ref{JC:prop:regular-disk-joint},
which controls the full disk scalar rather than the completed action.

More explicitly, cyclic equivariance diagonalizes every surviving quotient
mode $k$.  With the normalization
\begin{align*}
 \mathcal Q_{N,k}(d_k,w_k)
 &=q^{ff}_{N,k}|d_k|^2
   +2\operatorname {Re}(\overline{q^{fs}_{N,k}}d_k\overline{w_k})
   +q^{ss}_{N,k}|w_k|^2,\\
 \Gjp_{\alpha_*}(w_k)&=g_{N,k}|w_k|^2,
 \qquad g_{N,k}>0,
\end{align*}
write the polygonal half-Hessian on the same character as
\[
 \mathcal P_{N,k}(d_k,\zeta_k)
 =p^{ff}_{N,k}|d_k|^2
  +2\operatorname {Re}
    (\overline{p^{fs}_{N,k}}d_k\overline{\zeta_k})
  +p^{ss}_{N,k}|\zeta_k|^2,
 \qquad (Z_Nd)_k=Z_{N,k}d_k.
\]
Expanding the already proved polygonal reduction
\eqref{JC:eq:regular-action-polygon-reduction} gives the exact coefficients
\begin{equation}\label{JC:eq:regular-action-polygon-mode-coefficients}
 \boxed{
 \begin{aligned}
 q^{ss}_{N,k}&=p^{ss}_{N,k}-\frac12g_{N,k},\\
 q^{fs}_{N,k}&=p^{fs}_{N,k}+p^{ss}_{N,k}\overline{Z_{N,k}},\\
 q^{ff}_{N,k}&=3\kappa a^2
 {}+2\operatorname {Re}(p^{fs}_{N,k}Z_{N,k})
 +(p^{ss}_{N,k}+g_{N,k})|Z_{N,k}|^2.
 \end{aligned}}
\end{equation}
Thus the pure polygonal fan coefficient $p^{ff}_{N,k}$ cancels.  Before
positive margins are inserted, the remaining Schur complement is
\[
 3\kappa a^2+\frac12g_{N,k}|Z_{N,k}|^2
 -\frac{|p^{fs}_{N,k}+\frac12g_{N,k}\overline{Z_{N,k}}|^2}
 {p^{ss}_{N,k}-\frac12g_{N,k}}.
\]
This is an algebraic simplification, not a positivity assertion: the
polygonal support coefficient and the fan--support covector remain to be
estimated.

The same cancellation remains exact after margins are inserted.  Put
\[
 A_{N,k}:=p^{fs}_{N,k}
 +\frac12g_{N,k}\overline{Z_{N,k}},
 \qquad
 r_{N,k}:=p^{ss}_{N,k}
 -\left(\frac12+\mu_s\right)g_{N,k}.
\]
Here $A_{N,k}$ is precisely the regular fan derivative of the completed
covector $b_\alpha=\ell_\alpha-p_\alpha/2$ in the same orthonormal
Fourier normalization.  When $r_{N,k}>0$, the full modal Schur row is
equivalent to
\begin{equation}\label{JC:eq:regular-action-reduced-Schur}
 \boxed{
 (3\kappa-\mu_f)a^2
 +\left(\frac12-\mu_s\right)g_{N,k}|Z_{N,k}|^2
 -\frac{|A_{N,k}+\mu_sg_{N,k}\overline{Z_{N,k}}|^2}
 {r_{N,k}}\ge0.}
\end{equation}
Thus the mixed analytic blocker is a sharp all-character estimate for
$D_\alpha(\ell_\alpha-p_\alpha/2)$, not for the discarded covector
$\ell_\alpha-p_\alpha$.

Equation~\eqref{JC:eq:regular-action-Schur} is equivalent, mode by mode, to
\begin{equation}\label{JC:eq:regular-action-mode-matrix}
 \boxed{
 \begin{pmatrix}
  q^{ff}_{N,k}-\mu_fa^2&q^{fs}_{N,k}\\
  \overline{q^{fs}_{N,k}}&q^{ss}_{N,k}-\mu_sg_{N,k}
 \end{pmatrix}\succeq0.}
\end{equation}
Equivalently, with the usual quotient interpretation when the lower-right
entry vanishes,
\begin{align}
 q^{ss}_{N,k}-\mu_sg_{N,k}&\ge0,
 \label{JC:eq:regular-action-support-row}\\
 q^{ff}_{N,k}-\mu_fa^2
 &\ge\frac{|q^{fs}_{N,k}|^2}
 {q^{ss}_{N,k}-\mu_sg_{N,k}}.
 \label{JC:eq:regular-action-fan-Schur-row}
\end{align}
In particular, the pure support specialization already asks for
\begin{equation}\label{JC:eq:regular-action-support-gap}
 D_h^2\Phi_{\alpha_*}(h_*\mathbf1)[z,z]
 \ge(1+2\mu_s)\Gjp_{\alpha_*}(z).
\end{equation}
The exact fixed-fan support Hessian identifies the left side with the
regular-polygon reduced Dirichlet-to-Neumann block plus its geometric
area row.  More precisely, for the orthonormal character
$z_i=N^{-1/2}e^{\mathrm i k a i}$ let
$f_{N,k}=\mathcal I_{\alpha_*}z$,
$g_{N,k}=q_D(f_{N,k})$, let $p_N=-\partial_\nu u_N$, and let
$\psi_{N,k}$ be the character-$k$ reduced Helmholtz solution on the
area-$\pi$ regular polygon with boundary value $z_ip_N$ on side $i$.
Writing
\[
 \mathscr D_{N,k}=\int_{R_N}|\nabla\psi_{N,k}|^2
 -\lambda_N\int_{R_N}|\psi_{N,k}|^2,
\]
the exact support Hessian and area symbol give
\begin{equation}\label{JC:eq:regular-support-DtN-sector}
 \boxed{
 p^{ss}_{N,k}=\mathscr D_{N,k}
 +\frac{\lambda_N}{\pi}
   \frac{\cos(ka)-\cos a}{\sin a},
 \qquad
 q^{ss}_{N,k}=p^{ss}_{N,k}-\frac12g_{N,k}.}
\end{equation}
Hence the pure-support row is exactly the uniform sector inequality
\[
 \mathscr D_{N,k}
 +\frac{\lambda_N}{\pi}
   \frac{\cos(ka)-\cos a}{\sin a}
 \ge\left(\frac12+\mu_s\right)g_{N,k},
 \qquad 2\le k\le N-2.
\]
All terms except the regular-polygon reduced Helmholtz DtN energy are
explicit.  No closed terminal/LCA theorem proves this inequality, even
before the mixed Schur row
\eqref{JC:eq:regular-action-fan-Schur-row} is considered.

The exact Schwarz--Christoffel representation also shows why the DtN and
area rows must be estimated together.  After a rotation,
\[
 F_N'(z)=c_N(1-z^N)^{-2/N},\qquad J_N=|F_N'|,
\]
and, for the unnormalized character $z_i=e^{ikai}$, let
$\sigma_{N,k}$ be the step function on $\mathbb T$ which equals
$e^{ikai}$ on the preimage arc of side $i$ (its endpoint values are
irrelevant).  The pulled-back boundary trace is
\[
 h_{N,k}=\sigma_{N,k}c_N^{-1}
 |1-e^{iN\theta}|^{2/N}q_N,
 \qquad q_N=-\partial_r(u_N\circ F_N)|_{\mathbb T}.
\]
Uniform disk $W^{2,p}$ estimates give $q_N\to d_\circ$ uniformly,
including at the prevertices.  Keeping only the cross-prevertex
rectangles in the Gagliardo norm gives the harmonic-extension lower bound
$H(h_{N,k})\ge cN^2\{1-\cos(ka)\}-C$.  To pass to the reduced Helmholtz
energy, use the uniform boundary Friedrichs inequality
\[
 \|\phi\|_{L^2(R_N)}^2
 \le\varepsilon\|\nabla\phi\|_{L^2(R_N)}^2
 +C_\varepsilon\|h_{N,k}\|_{L^2(\partial R_N)}^2.
\]
Choose $\varepsilon<(2\sup_N\lambda_N)^{-1}$.  Rellich's identity gives
$\|h_{N,k}\|_{L^2(\partial R_N)}^2
=\int_{\partial R_N}p_N^2=2\lambda_N/h_N\le C$ in the unnormalized
character.  Hence every admissible reduced competitor has energy at least
$\frac12H(h_{N,k})-C$; the ground-state orthogonality only decreases the
competitor set.  This proves the rigorous DtN-only lower bound
\begin{equation}\label{JC:eq:corner-layer-DtN-lower}
 \mathscr D^{\rm un}_{N,k}
 \ge cN^2\{1-\cos(ka)\}-C.
\end{equation}
This does \emph{not} close the support row.  Formula
\eqref{JC:eq:regular-support-DtN-sector} uses the orthonormal character;
in the unnormalized convention its area term must also be multiplied by
$N$ and is of order
$-N^2\{1-\cos(ka)\}$.  The two leading corner-scale terms cancel in the
disk limit.  More precisely, the uncancelled coefficient contains
\[
 \frac{2}{a^2}
 \{\pi c_N^{-2}q_{N,*}^2-\lambda_N\}
 \{1-\cos(ka)\},
\]
for the prevertex endpoint amplitude $q_{N,*}$.  Rate-free convergence
alone does not control this coefficient after multiplication by $a^{-2}$.
One must consequently estimate the jointly renormalized remainder,
rather than combine
\eqref{JC:eq:corner-layer-DtN-lower} with an area symbol written in the other
normalization.  The next proposition supplies precisely that expansion,
uniformly through the mesoscopic regime $ka\downarrow0$.
\end{contract}

\begin{lemma}[Fixed-character polygon--disk support limit]
\label{JC:lem:fixed-character-support-limit}
For every fixed integer $k\ge2$,
\begin{equation}\label{JC:eq:fixed-character-support-limit}
 \frac{p^{ss}_{N,k}}{g_{N,k}}\longrightarrow1
 \qquad(N\to\infty).
\end{equation}
\end{lemma}

\begin{proof}
Use the unnormalized support character $z_i=e^{\mathrm i kia}$; the
normalization cancels from the quotient.  The material vertex velocity
at the vertex with bisector angle $\varphi_i=(i+\tfrac12)a$ is exactly
\[
 e^{\mathrm i k\varphi_i}\left\{
 \frac{\cos(ka/2)}{\cos(a/2)}e_r(\varphi_i)
 +\mathrm i\frac{\sin(ka/2)}{\sin(a/2)}
 e_\theta(\varphi_i)\right\}.
\]
For fixed $k$ its physical side-affine interpolant converges in
$H^{1/2}(\mathbb T;\mathbb C^2)$ to the smooth disk field
\[
 V_k^\circ(\theta)=e^{\mathrm i k\theta}
 \{e_r(\theta)+\mathrm i k e_\theta(\theta)\},
 \qquad V_k^\circ\cdot e_r=e^{\mathrm i k\theta}.
\]

We record the fixed-slot Hessian passage rather than citing an
unassembled stage.  Pull the area-normalized regular polygons to the
disk by the normalized near-identity conformal charts.  Their stiffness
forms and mass densities converge to the disk forms in every fixed
$L^p$, their first eigenfunctions converge strongly in $H_0^1$, and
the uniform convex-domain gradient bound makes
$\nabla u_N\otimes\nabla u_N\to
\nabla u_\circ\otimes\nabla u_\circ$ in $L^2$.
Polarize the exact volume formula for the normalized eigenvalue
Hessian with the fixed smooth extension of $V_k^\circ$ in both slots.
Every direct stiffness, mass, area, and product-of-first-jets row then
converges by the preceding strong coefficient convergence.

For the Schur row, put the fixed first-variation functional into the
reduced resolvent.  The spectral gap above the ground state is uniform,
the functionals converge in $H^{-1}$, and the resolvent identity on the
moving ground-state complements gives strong $H_0^1$ convergence of
the reduced states.  Thus the Schur pairing converges as well.  This
proves convergence of the complete polygonal normalized Hessian to the
disk Hessian on $V_k^\circ$; no mesh-ratio estimate is used because the
slot is fixed and smooth.

At the disk, scale and translation modes have already been projected
out and the normalized Hessian is
\[
 D^2\Phi_\circ[V_k^\circ,V_k^\circ]
 =2q_D(e^{\mathrm i k\theta}).
\]
On the other hand, the exact normal-sine interpolant
$f_{N,k}=\mathcal I_{\alpha_*}z$ converges to
$e^{\mathrm i k\theta}$ in $H^{1/2}$, so
$g_{N,k}\to q_D(e^{\mathrm i k\theta})$.
Since $p^{ss}_{N,k}$ is the polygonal half-Hessian coefficient,
\eqref{JC:eq:fixed-character-support-limit} follows.
\end{proof}

\begin{proposition}[Jointly renormalized regular support sector]
\label{JC:prop:regular-support-renormalized}
There is $N_s$ such that, for $N\ge N_s$ and every quotient character
$2\le k\le N-2$,
\begin{equation}\label{JC:eq:regular-support-closed-gap}
 \boxed{
 p^{ss}_{N,k}\ge\frac{11}{20}g_{N,k}
 =\left(\frac12+\frac1{20}\right)g_{N,k}.}
\end{equation}
Consequently the pure-support row
\eqref{JC:eq:regular-action-support-row} is closed with
$\mu_s=1/20$.
\end{proposition}

\begin{proof}
We keep one normalization until the last step.  By complex conjugation
assume $2\le k\le N/2$, and put
\[
 \beta=\frac2N,\qquad a=\pi\beta,\qquad
 x=\frac{k}{N}\in(0,\tfrac12],\qquad \xi=2\pi x.
\]
Use first the unnormalized complex character
$z_i=e^{2\pi\mathrm i x i}$.
Every quadratic value is then $N$ times its coefficient in the
orthonormal complex character.  The orthonormal real vector
$\sqrt{2/N}\operatorname{Re}z$ has the same modal coefficient.

Let
\[
 w_\beta(t)=|2\sin(\pi t)|^\beta,\qquad0<t<1.
\]
For the constant disk flux $d_\circ$, the exact corner trace is
$d_\circ\sigma_{N,k}w_\beta$.  Its only circular frequencies are
$k+mN$, and the beta integral gives
\begin{equation}\label{JC:eq:corner-Gamma-coefficient}
 I_{\beta,m}(x):=\int_0^1w_\beta(t)e^{-2\pi\mathrm i(m+x)t}\,dt
 =
 e^{-\pi\mathrm i(m+x)}
 \frac{\Gamma(1+\beta)}
 {\Gamma(1+\beta/2+m+x)\Gamma(1+\beta/2-m-x)}.
\end{equation}
Since $\pi d_\circ^2=\lambda_\circ$ by Rellich's identity, harmonic
energy and the area row, divided by $Nd_\circ^2$, combine into
\begin{equation}\label{JC:eq:corner-joint-finite-N-symbol}
 P_\beta(x):=
 2\pi\sum_{m\in\mathbb Z}|m+x||I_{\beta,m}(x)|^2
 +\frac{\cos(2\pi x)-\cos(\pi\beta)}{\sin(\pi\beta)}.
\end{equation}
The two terms in this formula separately contain the same
$\beta^{-1}$ pole with opposite signs.

We extract the finite part before estimating.  For $0<y<1$, put
\[
 T_\beta(y):=\sum_{n\ge0}(n+y)
 \left\{\frac{\Gamma(n+y-\beta/2)}
 {\Gamma(n+y+1+\beta/2)}\right\}^2.
\]
Reflection in \eqref{JC:eq:corner-Gamma-coefficient} gives the positive
identity
\begin{align}
 &\sum_m|m+x||I_{\beta,m}(x)|^2\notag\\
 &\quad=
 \frac{\Gamma(1+\beta)^2}{\pi^2}
 \left\{\sin^2\!\pi(x-\beta/2)T_\beta(x)
 +\sin^2\!\pi(x+\beta/2)T_\beta(1-x)\right\}.
 \label{JC:eq:corner-reflected-positive-series}
\end{align}
Separate the $n=0$ term.  Taylor's formula for $\log\Gamma$, together
with $|\psi_{\rm dg}^{(r)}(u)|\le C_r(1+u^{-r-1})$, gives, uniformly
for $0<\beta\le x/2$ and $x\le1/2$,
\begin{align*}
 T_\beta(x)
 &=\zeta(1+2\beta,x)
 +O\left(\frac\beta{x^2}
 +\frac{\beta|\log x|}{x}
 +\frac{\beta^2}{x^3}+\beta\right),\\
 T_\beta(1-x)&=\zeta(1+2\beta,1-x)+O(\beta),\\
 \zeta(1+2\beta,x)
 &=\frac1{2\beta}-\psi_{\rm dg}(x)
 +O\left(\frac{\beta|\log x|}{x}
 +\frac{\beta^2|\log x|^2}{x}+\beta\right),
\end{align*}
where $\psi_{\rm dg}=\Gamma'/\Gamma$.  Indeed, after comparison with
$(n+y)^{-1-2\beta}$ the derivative of the logarithmic quotient is
\[
 -2\psi_{\rm dg}(n+y)-(n+y)^{-1}+2\log(n+y)
 =O((n+y)^{-2}),
\]
so all $n\ge1$ remainders are absolutely summable; the singular powers
come only from $n=0$.

Expand the sine squares in
\eqref{JC:eq:corner-reflected-positive-series}, retaining their opposite
linear terms until cancellation, and use
$\Gamma(1+\beta)^2=1-2\gamma\beta+O(\beta^2)$.  We obtain
\begin{align}
 P_\beta(x)&=P(x)+
 O\left(\beta+\beta x|\log x|+\frac{\beta^2}{x}\right),
 \label{JC:eq:corner-mesoscopic-remainder}\\
 P(x)&:=-\frac{2\sin^2(\pi x)}{\pi}
 \{\psi_{\rm dg}(x)+\psi_{\rm dg}(1-x)+2\gamma\}.
 \label{JC:eq:corner-renormalized-symbol}
\end{align}

The high-character normal-sine disk symbol is
\begin{equation}\label{JC:eq:corner-disk-symbol}
 Q(x):=\frac{2\sin^4(\pi x)}{\pi^3}
 \{\zeta(3,x)+\zeta(3,1-x)\}.
\end{equation}
It satisfies $Q(x)\asymp x$ on $(0,1/2]$.  We now prove a strict
analytic gap.  Set
\[
 D(x):=-\{\psi_{\rm dg}(x)+\psi_{\rm dg}(1-x)+2\gamma\},
 \qquad H_3(x):=\zeta(3,x)+\zeta(3,1-x).
\]
The digamma and paired Hurwitz series give
\[
 D(x)=\frac1x+2\sum_{r\ge1}\zeta(2r+1)x^{2r},
\]
and
\[
 H_3(x)=x^{-3}+\sum_{n\ge1}\{(n-x)^{-3}+(n+x)^{-3}\}
 \le x^{-3}+14\zeta(3)-8.
\]
The tail increases with $x$ and the last constant is its value at
$x=1/2$.  Using $\sin(\pi x)\le\pi x$ and the integral-tail bound
$\zeta(3)<121/100$, we get
\begin{align*}
 D(x)-\frac35\frac{\sin^2(\pi x)}{\pi^2}H_3(x)
 &\ge\frac{2}{5x}+\frac{24-32\zeta(3)}5x^2\\
 &>\frac{2}{5x}-\frac{368}{125}x^2>0
 \qquad(0<x\le\tfrac12).
\end{align*}
Therefore
\begin{equation}\label{JC:eq:corner-three-fifths}
 \boxed{P(x)>\frac35Q(x),\qquad0<x\le\frac12.}
\end{equation}
Dividing \eqref{JC:eq:corner-mesoscopic-remainder} by $Q(x)\asymp x$
also shows
\begin{equation}\label{JC:eq:corner-relative-mesoscopic}
 P_\beta(x)=P(x)+o(Q(x))
\end{equation}
along every sequence for which $k=2x/\beta\to\infty$.

It remains to replace the constant flux by the true polygonal flux.
The exact area normalization of
$F_N'(z)=c_N(1-z^N)^{-\beta}$ is
\[
 c_N^{-2}=\sum_{m\ge0}
 \frac{(\beta)_m^2}{(m!)^2(mN+1)}.
\]
Since $(\beta)_m/m!\le C\beta m^{\beta-1}$ for $m\ge1$,
$c_N=1+O(\beta^3)$.  Put
$W_N=|F_N'|^2=c_N^2|1-z^N|^{-2\beta}$.  For every fixed $p>4$,
the change of variables $t=r^N$ gives
\begin{equation}\label{JC:eq:corner-weight-Lp-rate}
 \|W_N-1\|_{L^p(\mathbb D)}
 \le C_p\beta^{1+1/p}.
\end{equation}
Indeed, on $t\le1/2$ one has
$|\log|1-te^{\mathrm i\phi}||\le Ct$, while on $t\ge1/2$ all fixed
logarithmic moments remain uniformly integrable against
$|1-te^{\mathrm i\phi}|^{-2p\beta}$.  Since
$r\,dr=(\beta/2)t^{\beta-1}dt$, the $p$-th power of the norm is
$O(\beta^{p+1})$.

The pulled eigenfunction satisfies
\[
 -\Delta v_N=\lambda_NW_Nv_N,\qquad v_N|_{\mathbb T}=0,
 \qquad\int_{\mathbb D}W_Nv_N^2=1.
\]
The disk spectral gap first gives an $H^1$ estimate with the rate in
\eqref{JC:eq:corner-weight-Lp-rate}.  Apply a $W^{2,q}$ estimate with
$q>1$ to obtain $L^\infty$ control, substitute it back into the
difference equation, and then apply the $W^{2,p}$ estimate.  Thus
\begin{equation}\label{JC:eq:corner-flux-holder-rate}
 |\lambda_N-\lambda_\circ|
 +\|v_N-u_\circ\|_{W^{2,p}(\mathbb D)}
 \le C_p\beta^{1+1/p},
 \qquad
 \|q_N-d_\circ\|_{C^\alpha(\mathbb T)}
 \le C_p\beta^{1+1/p},
\end{equation}
where $q_N=-\partial_rv_N|_{\mathbb T}$ and
$\alpha=1-2/p>1/2$.

For $r\in C^\alpha$ and $f\in H^{1/2}(\mathbb T)$, the Gagliardo
double integral gives
\[
 [rf]_{H^{1/2}}^2
 \le2\|r\|_\infty^2[f]_{H^{1/2}}^2
 +C[r]_{C^\alpha}^2\|f\|_2^2.
\]
The positive Gamma series
\eqref{JC:eq:corner-reflected-positive-series} and
$T_\beta(y)\le C(\beta^{-1}+y^{-1})$ imply
\begin{equation}\label{JC:eq:corner-trace-energy-upper}
 H(\sigma_{N,k}w_\beta)\le Ck^2,
 \qquad2\le k\le N/2.
\end{equation}
Consequently polarization and
\eqref{JC:eq:corner-flux-holder-rate} give, with
$\varepsilon_N=C_p\beta^{1+1/p}$,
\begin{equation}\label{JC:eq:corner-true-flux-error}
 \left|H(\sigma_{N,k}c_N^{-1}w_\beta q_N)
 -d_\circ^2H(\sigma_{N,k}w_\beta)\right|
 \le C\varepsilon_Nk^2=o(k)
\end{equation}
uniformly along every growing-character sequence, because
$\varepsilon_Nk\le C\beta^{1/p}$.  The eigenvalue error in the area
row has the same bound.

The sector compatibility
$\int_{\partial R_N}z_ip_N\partial_\nu u_N\,ds=0$
follows from cyclic character orthogonality.  The uniform gap on
$u_N^\perp$, the Poisson $L^2$ estimate, and
$h_N\int_{\partial R_N}p_N^2\,ds=2\lambda_N$ show that reduced
Helmholtz energy differs from harmonic extension by $O(1)$ in the
unnormalized character.  This is also $o(k)$ when $k\to\infty$.
It follows from
\eqref{JC:eq:corner-relative-mesoscopic} and
\eqref{JC:eq:corner-true-flux-error} that
\begin{equation}\label{JC:eq:corner-polygon-symbol-limit}
 \frac{p^{ss,\mathrm{un}}_{N,k}}{Nd_\circ^2}
 =P(k/N)+o(Q(k/N))
\end{equation}
on every growing-character sequence.

Finally, the exact Bessel alias formula for $g_{N,k}$, together with
\[
 \omega_n=|n|+1+O(|n|^{-1}),\qquad
 (1-n^2)^{-2}=n^{-4}\{1+O(n^{-2})\},
\]
gives
\begin{equation}\label{JC:eq:corner-disk-alias-limit}
 \frac{g^{\mathrm{un}}_{N,k}}{Nd_\circ^2}
 =Q(k/N)\{1+o(1)\}
\end{equation}
uniformly for $k\to\infty$.  Indeed every alias has
$|k+mN|\ge k$, and replacing
$\cos(ka)-\cos a$ by $\cos(ka)-1$ costs $O(k^{-2})$
relatively.

If \eqref{JC:eq:regular-support-closed-gap} failed along a sequence, choose
a violating character and conjugate it into $2\le k_N\le N/2$.
If $k_N$ is bounded, pass to a constant subsequence; the retained
one-fixed-slot disk limit gives
$p^{ss}_{N,k_N}/g_{N,k_N}\to1$.  If $k_N\to\infty$, then
\eqref{JC:eq:corner-three-fifths},
\eqref{JC:eq:corner-polygon-symbol-limit}, and
\eqref{JC:eq:corner-disk-alias-limit} give
\[
 \liminf_N\frac{p^{ss}_{N,k_N}}{g_{N,k_N}}\ge\frac35.
\]
Both alternatives contradict a ratio below $11/20$.  This proves
\eqref{JC:eq:regular-support-closed-gap}.
\end{proof}

\begin{lemma}[Polygonal mixed Schwarz--Christoffel finite part]
\label{JC:lem:regular-polygonal-mixed-finite-part}
Fix $0<x<1/2$, put
\[
 \theta=\pi x,\qquad \rho=e^{2\pi ix},\qquad
 D(x)=-\{\psi(x)+\psi(1-x)+2\gamma\},
 \qquad T(x)=\psi_1(x)-\psi_1(1-x)=-D'(x).
\]
For $\beta=2/N$, $a=\pi\beta$, and a regular fan character with
$k/N\to x$, the constant-flux polygonal fan--support coefficient has the
phase-realified limit
\begin{equation}\label{JC:eq:regular-polygonal-mixed-finite-part}
 \boxed{
 L_{\rm poly}(x):=
 \lim_{\beta\downarrow0}
 \left[-e^{i\theta}\frac{p^{fs}_\beta}{d_\circ^2a}\right]
 =\frac{\cos\theta}{2\pi}D(x)
 -\frac{\sin\theta}{2\pi^2}T(x)
 =\frac{\sin\theta}{2\pi^2}
 \{D'(x)+\pi\cot\theta D(x)\}.}
\end{equation}
Let
\[
 F_p(y):=\sum_{n\ge0}\frac{(-1)^n}{(n+y)^p}
 =2^{-p}\{\zeta(p,y/2)-\zeta(p,(y+1)/2)\}
\]
and define the exact disk-forcing symbol
\begin{equation}\label{JC:eq:regular-disk-forcing-finite-part}
 R_I(x)=-\frac{\sin\theta}{4\pi^3}
 \{F_3(x)-F_3(1-x)\}
 +\frac{\cos\theta}{4\pi^2}
 \{F_2(x)+F_2(1-x)\}.
\end{equation}
Then the genuine polygonal completed-covector finite part exists and is
\begin{equation}\label{JC:eq:regular-completed-covector-finite-part}
 \boxed{
 R_{b,{\rm poly}}(x)=
 \frac{D'(x)+\pi\cot\theta D(x)}{2\pi^2}-R_I(x),
 \qquad
 \frac{A_{N,k}}{d_\circ^2a}
 =-e^{-i\theta}\sin\theta R_{b,{\rm poly}}(x)+o_x(1).}
\end{equation}
The formulas extend continuously to $x=1/2$, where both
$R_I$ and $R_{b,{\rm poly}}$ vanish.
\end{lemma}

\begin{proof}
On one regular Schwarz--Christoffel side let
\[
 w_\beta(t)=|2\sin\pi t|^\beta,\qquad
 J_\beta(t)=\int_{1/2}^t w_\beta(u)^{-1}\,du,\qquad 0<t<1.
\]
For the complex fan character $d_i=\rho^i$, differentiating the turning
angles gives $\dot\theta_i=(\rho-1)^{-1}\rho^i$ and $\dot h=0$.
Centred physical arclength satisfies
$ds=c_Naw_\beta^{-1}dt$ and $s=c_NaJ_\beta(t)$, while the pulled-back
flux is $p_N=c_N^{-1}w_\beta q_N$.  Hence the exact fan--flux datum is
\[
 -a(\rho-1)^{-1}\rho^i w_\beta(t)J_\beta(t)q_N(t).
\]
For $y=m+x$ put
\[
 I_\beta(y)=\int_0^1w_\beta(t)e^{-2\pi iyt}\,dt,\qquad
 S_\beta(y)=\int_0^1w_\beta(t)J_\beta(t)e^{-2\pi iyt},dt,
\]
and $M_\beta(x)=\sum_m|m+x|S_\beta(m+x)\overline{I_\beta(m+x)}$.
The harmonic mixed row is
$-2\pi a(\rho-1)^{-1}M_\beta(x)$ after division by $d_\circ^2$; the area
row is
\[
 \frac14h_N\sec^2(a/2)(1+\rho^{-1}),\qquad
 h_N=\sqrt{\frac\pi{N\tan(a/2)}}.
\]

It remains to extract the joint finite part.  Set
\[
 C_\beta=\int_0^{1/2}w_\beta^{-1}
 =2^{-\beta}\frac{\Gamma((1-\beta)/2)}
 {2\sqrt\pi\,\Gamma(1-\beta/2)}.
\]
Then $C_0=1/2$, $C'_0=0$, and direct integration at the two endpoints
gives
\begin{align*}
 w_\beta(t)J_\beta(t)
 &=-C_\beta(2\pi t)^\beta+\frac{t}{1-\beta}
   +O(t^{\beta+2}+t^3),\\
 w_\beta(1-u)J_\beta(1-u)
 &= C_\beta(2\pi u)^\beta-\frac{u}{1-\beta}
   +O(u^{\beta+2}+u^3).
\end{align*}
The endpoint expansion can be used here with a summable differentiated
remainder, as follows.  Choose smooth cutoffs at $0$ and $1$, subtract
from $w_\beta$ its two endpoint models of type $t^\beta$, and subtract
from $w_\beta J_\beta$ the four endpoint models displayed above.  The
remainders and their first $\beta$--derivatives have two integrable
derivatives; differentiating in $\beta$ merely inserts a logarithm.

Here is the resulting coefficient ledger, which also fixes the two tail
signs.  Put
\[
 a_\beta=\frac{\Gamma(1+\beta)}{2\pi},\qquad
 \phi_\beta=\frac\pi2(1+\beta).
\]
For $\sigma=\operatorname{sgn}y$, the two endpoint models give
\begin{align*}
 A_+(\beta,x)&=a_\beta
  \{e^{-i\phi_\beta}+e^{-2\pi ix}e^{i\phi_\beta}\},&
 B_+(\beta,x)&=C_\beta a_\beta
  \{-e^{-i\phi_\beta}+e^{-2\pi ix}e^{i\phi_\beta}\},\\
 A_-(\beta,x)&=a_\beta
  \{e^{i\phi_\beta}+e^{-2\pi ix}e^{-i\phi_\beta}\},&
 B_-(\beta,x)&=C_\beta a_\beta
  \{-e^{i\phi_\beta}+e^{-2\pi ix}e^{-i\phi_\beta}\}.
\end{align*}
Thus, for $j=0,1$, fixed $x\in(0,1/2]$, and
$0\le\beta\le\beta_x<1$,
\begin{align}
 I_\beta(y)&=A_\sigma(\beta,x)|y|^{-1-\beta}
             +E^I_\beta(y),\notag\\
 S_\beta(y)&=B_\sigma(\beta,x)|y|^{-1-\beta}
             +E^S_\beta(y),\notag\\
 |\partial_\beta^jE^I_\beta(y)|
 +|\partial_\beta^jE^S_\beta(y)|
 &\le C_x\frac{\{1+\log(2+|y|)\}^{j}}{(1+|y|)^2}.
 \label{JC:eq:regular-mixed-Fourier-remainders}
\end{align}
The linear endpoint models are included in $E^S$; their coefficients are
$O(|y|^{-2})$.  The displayed estimate follows by two integrations by
parts on the cutoff remainders.  It is uniform down to $\beta=0$ because
$t^\beta\log t$ and the second derivatives of
$t^{\beta+2}\log t$ are integrable uniformly on
$0\le\beta\le\beta_x$.

Direct multiplication gives
\[
 B_+(\beta,x)\overline{A_+(\beta,x)}
 =2iC_\beta a_\beta^2\sin(2\pi x-\pi\beta),
 \qquad
 B_-(\beta,x)\overline{A_-(\beta,x)}
 =2iC_\beta a_\beta^2\sin(2\pi x+\pi\beta).
\]
If $c_\pm(\beta,x)$ denote these two coefficients and
\[
 R_m(\beta,x):=|m+x|S_\beta(m+x)\overline{I_\beta(m+x)}
 -c_{\operatorname{sgn}(m+x)}(\beta,x)|m+x|^{-1-2\beta},
\]
then~\eqref{JC:eq:regular-mixed-Fourier-remainders}, followed by the product
rule, gives
\begin{equation}\label{JC:eq:regular-mixed-endpoint-remainder}
 |R_m(\beta,x)|+|\partial_\beta R_m(\beta,x)|
 \le C_x\frac{1+\log(2+|m|)}{(1+|m|)^2}.
\end{equation}
Indeed the coefficients of every subtracted singular model follow from
\[
 \int_0^\infty t^\alpha e^{-2\pi iyt}\,dt
 =\Gamma(1+\alpha)e^{-i\pi(1+\alpha)/2}
  (2\pi y)^{-1-\alpha},\qquad y>0,
\]
first with an exponential cutoff and then by continuation.  This proves
the displayed bound rather than assuming a formal asymptotic series.
Its majorant is summable, so the subtracted series is $C^1$ at
$\beta=0$ and termwise passage to the finite part is justified.  The
only non-summable tail coefficients are
\[
 c_\pm(\beta,x)=2iC_\beta
 \left\{\frac{\Gamma(1+\beta)}{2\pi}\right\}^{\!2}
 \sin(2\pi x\mp\pi\beta),
\]
multiplying the corresponding $|m+x|^{-1-2\beta}$ tails.  The linear
endpoint terms produce summable order $|m|^{-2-\beta}$ terms, so no
additional $\beta\zeta(1+\beta)$ finite part is possible.

More explicitly,
\[
 M_\beta(x)=c_+(\beta,x)\zeta(1+2\beta,x)
 +c_-(\beta,x)\zeta(1+2\beta,1-x)
 +\sum_{m\in\mathbb Z}R_m(\beta,x).
\]
At $\beta=0$ the exact formulas below give
\[
 R_m(0,x)=-\frac{i\sin^2\theta}{2\pi^3}
 \frac{\operatorname{sgn}(m+x)}{|m+x|^2},
 \qquad
 \sum_mR_m(0,x)=-\frac{i\sin^2\theta}{2\pi^3}T(x).
\]
Also $C'_0=0$ and
$\partial_\beta\Gamma(1+\beta)^2|_{\beta=0}=-2\gamma$; the opposite
$\mp\pi\cos(2\theta)$ derivatives of the two sine coefficients cancel.
These identities account for every finite term in the next display.

At $\beta=0$ one has
\[
 I_0(y)=e^{-\pi iy}\frac{\sin\pi y}{\pi y},\qquad
 S_0(y)=\frac i2I_0(y)
 \left(\cot\theta-\frac1{\pi y}\right).
\]
Using
$\zeta(1+2\beta,q)=(2\beta)^{-1}-\psi(q)+O(\beta)$ in the two tails and
summing the $C^1$ remainder from
\eqref{JC:eq:regular-mixed-endpoint-remainder} yields, locally uniformly in
$x$,
\[
 M_\beta(x)=\frac{i\sin(2\theta)}{4\pi^2\beta}
 +i\left\{\frac{\sin(2\theta)}{4\pi^2}D(x)
 -\frac{\sin^2\theta}{2\pi^3}T(x)\right\}+O_x(\beta).
\]
Because $2i(\rho-1)^{-1}\sin(2\theta)=1+\rho^{-1}$ and
$h_N\sec^2(a/2)=1+O(a^2)$, the pole cancels the area row.  Multiplication
by $-e^{i\theta}$ proves~\eqref{JC:eq:regular-polygonal-mixed-finite-part}.
The exact disk alias calculation gives
\eqref{JC:eq:regular-disk-forcing-finite-part}; adding the corrector
contribution $-\sin\theta R_I$ proves
\eqref{JC:eq:regular-completed-covector-finite-part}.  Reflection about
$x=1/2$ gives the stated continuous endpoint values.
\end{proof}

\begin{proposition}[Strict support--double-disk gap]
\label{JC:prop:regular-support-double-disk-gap}
For every $0<x\le1/2$ one has
\begin{equation}\label{JC:eq:regular-support-double-disk-gap}
 \boxed{P(x)>2Q_{\mathrm{mat}}(x)},\qquad
 Q_{\mathrm{mat}}(x):=\frac{\sin^4(\pi x)}{\pi^3}
 \{\zeta(3,x)+\zeta(3,1-x)\}.
\end{equation}
\end{proposition}

\begin{proof}
Write $H_3(x)=\zeta(3,x)+\zeta(3,1-x)$.  The alternating Taylor
series for $\sin^2y$ on $0\le y\le\pi/2$ gives
\[
 \sin^2(\pi x)\le \pi^2x^2(1-2x^2).
\]
Indeed the alternating expansion gives
$\sin^2y\le y^2-y^4/3+2y^6/45$.  Since $y\le\pi/2$ and
$1/3-2/\pi^2-\pi^2/90>0$, its right side is at most
$y^2-(2/\pi^2)y^4$.  The convergent series for $D$ gives
\[
 D(x)=\frac1x+2x^2\sum_{n\ge1}\frac1{n(n^2-x^2)}
 \ge\frac1x+2\zeta(3)x^2.
\]
Moreover $H_3(x)-x^{-3}$ is increasing on $[0,1/2]$, term by term, and
therefore
\[
 H_3(x)\le x^{-3}+H_3(1/2)-8
 =x^{-3}+14\zeta(3)-8.
\]
Using the elementary bounds $6/5<\zeta(3)<121/100$ (sum the first ten
terms and bound the decreasing tail between
$\int_{11}^\infty t^{-3}dt$ and $\int_{10}^\infty t^{-3}dt$), it follows that
\begin{align*}
 D(x)-\frac{\sin^2(\pi x)}{\pi^2}H_3(x)
 &\ge 2x+(8-12\zeta(3))x^2+(28\zeta(3)-16)x^4\\
 &\ge x\left(2-\frac{163}{25}x+\frac{88}{5}x^3\right)>0.
\end{align*}
For the last inequality, the cubic bracket has its minimum at
$x_0=\sqrt{163/1320}<44/125$, where it equals
$2-(326/75)x_0>0$.  Since
$P(x)>2Q_{\mathrm{mat}}(x)$ is exactly the last displayed strict
inequality multiplied by $2\sin^2(\pi x)/\pi$, the proof is complete.
\end{proof}

\begin{remark}[Retraction of the old period-three completion]
\label{JC:rem:period-three-completion-retracted}
The revision-21 period-three completion used the material support constant
$E_3$.  The jointly renormalized Schwarz--Christoffel calculation above
shows instead that the true real-paired polygonal support finite part is
\[
 \frac12P(1/3)=\frac{9\log3}{4\pi}=0.786823093\ldots,
\]
whereas the material computation gives $E_3=0.77956467\ldots$.  The
missing difference is a same-order corner finite part.  Therefore the old
period-three proposition remains withdrawn.  The proofs below use only
the genuine polygonal diagonal and mixed finite parts computed above; they
do not reinstate the material transfer.
\end{remark}

\begin{proposition}[Certified completed-covector ratio]
\label{JC:prop:regular-completed-covector-ratio}
For $0<x<1/2$, put
\[
 J(x):=\frac{D'(x)+\pi\cot(\pi x)D(x)}{2\pi^2}.
\]
Then
\begin{equation}\label{JC:eq:regular-completed-covector-ratio}
 \boxed{2R_I(x)<J(x)<R_I(x)<0.}
\end{equation}
Consequently, for $0\le\mu\le1/4$ and
$Q_{\rm mat}=Q/2$,
put
\[
 D_{\rm poly}(x;\mu)
 :=P(x)-(1+2\mu)Q_{\rm mat}(x).
\]
Then $D_{\rm poly}(x;\mu)>0$ and
\begin{equation}\label{JC:eq:regular-completed-covector-determinant}
 \begin{aligned}
 \Delta_\mu(x):={}&(1-2\mu)R_I(x)^2
 D_{\rm poly}(x;\mu)\\
 &-Q_{\rm mat}(x)
 \{R_{b,{\rm poly}}(x)-2\mu R_I(x)\}^2>0
 \qquad(0<x<1/2).
 \end{aligned}
\end{equation}
At $x=1/2$ both covectors vanish and $\Delta_\mu(1/2)=0$.
\end{proposition}

\begin{proof}
Write $s=\sin(\pi x)$, $c=\cos(\pi x)$ and
$A=F_2(x)+F_2(1-x)$.  Since
$A'=-2\{F_3(x)-F_3(1-x)\}$, one has the exact identities
\begin{equation}\label{JC:eq:regular-ratio-derivative-identities}
 J=\frac{(sD)'}{2\pi^2s},\qquad
 R_I=\frac{(s^2A)'}{8\pi^3s}.
\end{equation}
Thus it suffices to prove the three strict signs
\[
 -R_I>0,\qquad U:=R_I-J>0,\qquad V:=J-2R_I>0.
\]
We give the finite rational certificate, including the endpoint
regularizations, so that no sampled floating-point value is used.

At the left endpoint the convergent expansions are
\begin{align*}
 D&=x^{-1}+2\sum_{r\ge1}\zeta(2r+1)x^{2r},\\
 \pi\cot(\pi x)&=x^{-1}-2\sum_{r\ge1}\zeta(2r)x^{2r-1},\\
 F_3(x)-F_3(1-x)
 &=x^{-3}-2\sum_{r\ge0}\binom{2r+2}{2}\eta(2r+3)x^{2r},\\
 F_2(x)+F_2(1-x)
 &=x^{-2}+4\sum_{r\ge0}(r+1)\eta(2r+3)x^{2r+1}.
\end{align*}
The singular terms cancel symbolically and give
\[
 J(0)=-\frac16,\qquad R_I(0)=-\frac1{12},\qquad
 \lim_{x\downarrow0}\frac{V(x)}x
 =\frac{3\zeta(3)}{4\pi^2}>0.
\]
Ten retained orders, $q=x^2$, and the elementary bounds
$\zeta(r)<121/100$ for $r\ge3$ and $\zeta(r)<5/3$ for $r\ge2$
bound the omitted $D,D',D/x,\cot$, and divided-cotangent tails by
\begin{gather*}
 \frac{121}{50}\frac{x^{22}}{1-q},\quad
 \frac{121}{25}x^{21}\frac{11-10q}{(1-q)^2},\quad
 \frac{121}{50}\frac{x^{21}}{1-q},\\
 \frac{10}{3}\frac{x^{21}}{1-q},\qquad
 \frac{10}{3}\frac{x^{20}}{1-q}.
\end{gather*}
The two alternating-$F$ tails are dominated by
\[
 \frac{8(12)^2q^{11}}{(1-q)^3},\qquad
 \frac{4x(12)q^{11}}{(1-q)^2},
\]
with the divided second bound obtained by deleting $x$.  Directed rational
evaluation on $0\le x\le1/8$ then gives
\begin{equation}\label{JC:eq:regular-ratio-left-certificate}
 R_I<-0.064086,\qquad U>0.035346,\qquad V/x>0.034494.
\end{equation}

For the right endpoint put $y=1/2-x$ and $C=\cos(\pi y)$.  Then
\begin{align*}
 D&=4\log2+2\sum_{n\ge1}(2^{2n+1}-1)\zeta(2n+1)y^{2n},\\
 A&=\sum_{n\ge0}(2n+1)2^{2n+3}\beta_D(2n+2)y^{2n}.
\end{align*}
If $\mathcal A=(CD)_y$ and $\mathcal B=(C^2A)_y$, the endpoint zeros
factor exactly as
\[
 -R_I/y=\frac{\mathcal B/y}{8\pi^3C},\quad
 U/y=\frac{4\pi\mathcal A/y-\mathcal B/y}{8\pi^3C},\quad
 V/y=\frac{\mathcal B/y-2\pi\mathcal A/y}{4\pi^3C}.
\]
With nine retained orders and $q=4y^2$, the positive tails for
$D,D_y/y,A,A_y/y$ are bounded respectively by
\[
 \frac{121}{25}\frac{q^{10}}{1-q},\quad
 \frac{968}{25}q^9\frac{10-9q}{(1-q)^2},\quad
 \frac{8(21)q^{10}}{(1-q)^2},\quad
 \frac{128(10)(21)q^9}{(1-q)^3}.
\]
Directed rational evaluation on $0\le y\le1/32$ gives
\begin{equation}\label{JC:eq:regular-ratio-right-certificate}
 -R_I/y>0.167528,\qquad U/y>0.105954,\qquad V/y>0.001213.
\end{equation}

It remains only the compact interval $[1/8,15/32]$.  The executable
certificate
\texttt{tests/check\_jc\_ar\_ratio\_symbol.py} partitions it into the
11,264 exact rational cells
\[
 [i/32768,(i+1)/32768],\qquad4096\le i\le15359.
\]
On each cell it evaluates twelve terms of the positive $D,D'$ series and
the alternating $F_2,F_3$ series.  Their omitted tails are bounded by
\[
 \frac{4x^2}{3\cdot12^2},\qquad
 \frac{32x}{9\cdot12^2},\qquad (12+z)^{-p},
\]
respectively.  The rational enclosure
\[
 \frac{103993}{33102}<\pi<\frac{104348}{33215}
\]
comes from Machin's identity; sine and cosine use alternating Taylor
series with their signed next terms.  Every arithmetic operation is
rounded outwards at 64 decimal digits.  In particular, sign reversal is
exact and integer powers use only directed interval multiplication.  The
finite minima are
\begin{equation}\label{JC:eq:regular-ratio-middle-certificate}
 -R_I>0.0056578993130883869,\quad
 U>0.0038837669640606282,\quad
 V>0.0009703006655867654.
\end{equation}
Equations~\eqref{JC:eq:regular-ratio-left-certificate}--
\eqref{JC:eq:regular-ratio-middle-certificate} prove
\eqref{JC:eq:regular-completed-covector-ratio}.  The script was independently
audited against exact rational arithmetic, including a hostile seven-digit
ambient decimal context; the same 64-digit enclosures result.

Finally Lemma~\ref{JC:lem:regular-polygonal-mixed-finite-part} gives
$R_{b,{\rm poly}}=J-R_I$.  Hence
$0<t:=R_{b,{\rm poly}}/R_I<1$.  For $0\le\mu\le1/4$,
$|t-2\mu|\le1-2\mu$.  Therefore
\[
 \Delta_\mu\ge(1-2\mu)R_I^2\{P-2Q_{\rm mat}\}>0
\]
by Proposition~\ref{JC:prop:regular-support-double-disk-gap}.  This proves
\eqref{JC:eq:regular-completed-covector-determinant}.
\end{proof}

\begin{lemma}[Uniform polygonal mixed bridge]
\label{JC:lem:regular-mixed-uniform-bridge}
Let $\beta=2/N$, $a=\pi\beta$, $x=k/N$, and
$2\le k\le N/2$.  The constant-flux harmonic-plus-area mixed coefficient
$p^{fs,0}_{\beta}$ satisfies
\begin{equation}\label{JC:eq:regular-mixed-CF2P}
 \left|-e^{i\pi x}\frac{p^{fs,0}_{\beta}}{d_\circ^2a}
       -L_{\rm poly}(x)\right|
 \le C\left(\beta+\frac{\beta^2}{x}\right).
\end{equation}
The actual reduced Helmholtz-minus-harmonic mixed modal row satisfies
\begin{equation}\label{JC:eq:regular-mixed-RR-ABS}
 \boxed{|E^{fs}_{N,k}|\le C_pa^2.}
\end{equation}
After replacing the constant flux by the true polygonal flux, the complete
mixed and diagonal ledgers have the following consequence.  If $k_N\to
\infty$, their Schur quotient differs from the limiting proxy determined
by~\eqref{JC:eq:regular-completed-covector-determinant} by $o(a^2)$; if
$k_N$ is bounded, the entire mixed Schur quotient is $o(a^2)$.
\end{lemma}

\begin{proof}
For~\eqref{JC:eq:regular-mixed-CF2P}, keep the $m=0$ term of
$M_\beta(x)$ exact.  Reflection gives $S_\beta(0)=0$, so that term is
$x^2O_{C^2}(1)$.  On $|m|\ge1$, subtract the $t^\beta$ endpoint models
from $I_{\beta,m}$ and the $t^\beta$ and linear endpoint models from
$S_{\beta,m}$.  Two integrations by parts give, for
$0\le\ell\le1$ and $0\le j\le2$,
\begin{equation}\label{JC:eq:regular-mixed-joint-remainder}
 |\partial_x^\ell\partial_\beta^jR_m(\beta,x)|
 \le C\frac{1+\log^2(2+|m|)}{(1+|m|)^2}.
\end{equation}
The removed $|m+x|^{-1-2\beta}$ families are first summed as paired
Hurwitz zeta functions.  After subtracting their common pole and finite
part, their Laurent remainders are jointly $C^{1,2}$ for
$0\le\beta\le x\le1/2$.  Reflection makes the first $\beta$ coefficient
vanish at $x=0$.  Two-variable Taylor's formula therefore gives
\[
 \left|M_\beta(x)-\frac{i\sin(2\pi x)}{4\pi^2\beta}
 -i\left\{\frac{\sin(2\pi x)}{4\pi^2}D(x)
 -\frac{\sin^2(\pi x)}{2\pi^3}T(x)\right\}\right|
 \le C(\beta x+\beta^2).
\]
The exact area pole cancellation proves~\eqref{JC:eq:regular-mixed-CF2P}.

We prove~\eqref{JC:eq:regular-mixed-RR-ABS}, which is the reduced row omitted
in the first bridge attempt.  Pull the regular polygon to the disk and
write $V_N=\lambda_N|F_N'|^2$.  In a nontrivial character sector let
$U_Nf$ solve $(-\Delta-V_N)U_Nf=0$ with boundary value $f$, and set
\[
 \mathcal K_N(f,g)=\langle(\Lambda_{V_N}-\Lambda_0)f,g\rangle.
\]
The ground state is invariant, so every retained sector is orthogonal to
it.  The first/second spectral gap is uniform; hence, for fixed $p>4$ and
$q=2p/(p-1)$,
\begin{equation}\label{JC:eq:regular-sector-solution-bound}
 \|U_Nf\|_{L^q}+\|Hf\|_{L^q}\le C_p\|f\|_\infty.
\end{equation}
The retained SC weight estimate gives
$\|V_N-\lambda_\circ\|_{L^p}\le C_p\beta^{1+1/p}$.
Alessandrini's identity and~\eqref{JC:eq:regular-sector-solution-bound} yield
\begin{equation}\label{JC:eq:regular-sector-potential-perturbation}
 |\mathcal K_N(f,g)-\mathcal K_\circ(f,g)|
 \le C_p\beta^{1+1/p}\|f\|_\infty\|g\|_\infty.
\end{equation}

For the unnormalized constant-flux traces,
\[
 f^s=d_\circ\sigma_{N,k}w_\beta,\qquad
 f^f=-d_\circ a(\rho-1)^{-1}\sigma_{N,k}w_\beta J_\beta.
\]
Their Fourier coefficients at $n=k+mN$ are, apart from the displayed
fixed factors, $I_{\beta,m}$ and $S_{\beta,m}$.  Uniformly for
$0<\beta\le x\le1/2$,
\[
 |I_{\beta,0}|\le C,\quad |S_{\beta,0}|\le Cx,\qquad
 |I_{\beta,m}|\le\frac{Cx}{1+|m|},\quad
 |S_{\beta,m}|\le\frac C{1+|m|}\quad(m\ne0).
\]
For the constant disk potential the reduced-minus-harmonic multiplier is
\[
 \varkappa_n=j_{0,1}\frac{J_{|n|}'(j_{0,1})}{J_{|n|}(j_{0,1})}-|n|
 =-j_{0,1}\frac{J_{|n|+1}(j_{0,1})}{J_{|n|}(j_{0,1})},
 \qquad |\varkappa_n|\le\frac C{1+|n|}.
\]
Here $|n|\ge2$, so there is no ground-state resonance.  It follows that
\[
 |\mathcal K_\circ(f^f,f^s)|
 \le C\frac ax\left\{\frac{x}{1+Nx}
 +\sum_{m\ne0}\frac{x}{N(1+|m|)^3}\right\}\le Ca.
\]
The true flux obeys
$\|q_N-d_\circ\|_\infty\le C_pa^{1+1/p}$; the Green identity
$\mathcal K_N(f,g)=-\int V_NU_Nf\,\overline{Hg}$ and
\eqref{JC:eq:regular-sector-solution-bound} show that replacing either slot
changes the preceding unnormalized form by $o(a)$.  Finally an
unnormalized character is $\sqrt N$ times an orthonormal one in each
mixed slot.  Division by $N=2\pi/a$ proves
\eqref{JC:eq:regular-mixed-RR-ABS}.

For completeness, define the positive side-mass normalization by
\[
 \bar q_N^2:=
 \frac{\int_0^1w_\beta q_N^2}{\int_0^1w_\beta}.
\]
The regular-side stationarity identity
$\int_{S_i}p_N^2\,ds=(\lambda_N/\pi)L_i$, together with
$p_N=c_N^{-1}w_\beta q_N$ and
$ds=c_Na w_\beta^{-1}dt$, gives exactly
\[
 \bar q_N^2=\frac{\lambda_N}{\pi}c_N^2
 \frac{\int_0^1w_\beta^{-1}}{\int_0^1w_\beta}
 =\frac{\lambda_N}{\pi}\{1+O(\beta^3)\}.
\]
Here $c_N=1+O(\beta^3)$ and the expansions of
$\int w_\beta^{\pm1}$ have the same quadratic coefficient because
$\int_0^1\log|2\sin\pi t|\,dt=0$.  Moreover $q_N$ is cell-periodic and
$\bar q_N$ lies between its minimum and maximum on a cell.  Rescaling a
cell of length $O(\beta)$ to $[0,1]$ in
\eqref{JC:eq:corner-flux-holder-rate} therefore proves
\[
 \|q_N-\bar q_N\|_{C^\alpha_{\rm cell}}
 \le C_p\beta^{2-1/p},\quad \alpha=1-2/p>1/2.
\]
The endpoint coefficient bounds used above also give
\[
 \sum_m|m+x||I_{\beta,m}|^2\le C\frac{x^2}{\beta},\qquad
 \sum_m|m+x||S_{\beta,m}|^2\le\frac C\beta.
\]
Multiplication by a $C^\alpha_{\rm cell}$ function is bounded on this
Bloch $H^{1/2}$ space; polarization and
$a|\rho-1|^{-1}\le C\beta/x$ therefore show that replacing the constant
flux by its cell oscillation changes the harmonic mixed row by at most
$C_p\beta^{2-1/p}$.  The disk aliases give
uniformly
\begin{equation}\label{JC:eq:regular-mixed-alias-ledger}
 cx\le g_{N,k}\le Cx,\qquad |Z_{N,k}|\le Ca,\qquad
 |g_{N,k}Z_{N,k}|\le Cax.
\end{equation}
On every growing-character sequence,~\eqref{JC:eq:regular-mixed-CF2P},
\eqref{JC:eq:regular-mixed-RR-ABS}, the support asymptotic, and the preceding
true-flux estimate give, for the limiting proxies denoted by bars,
\begin{align*}
 |\bar A|&\le Cax,& r,\bar r&\ge cx,\\
 |A-\bar A|&\le o(ax)+C_pa^{2-1/p}+Ca^2,&
 |r-\bar r|&=o(x).
\end{align*}
Elementary perturbation of $|A+\mu_sg\overline Z|^2/r$ then costs
\[
 o(a^2x)+C_pa^{3-1/p}+C_pa^{4-2/p}/x+O(a^3+a^4/x)=o(a^2),
\]
because $x\ge\beta\asymp a$ and $p>4$.  If $k$ stays bounded, the
fixed-slot limit and~\eqref{JC:eq:regular-mixed-alias-ledger} give
$r\ge c_ka$, while the numerator is
$O_k(a^2)+O_p(a^{2-1/p})$; its quotient is
$O_p(a^{3-2/p})=o(a^2)$.  This proves the last assertion.
\end{proof}

\begin{proposition}[Closure of the regular completed-action Schur block]
\label{JC:prop:regular-action-Schur-closed}
For all sufficiently large $N$, Contract~\ref{JC:con:JC-AR} holds with
\begin{equation}\label{JC:eq:regular-action-closed-constants}
 \boxed{\mu_s=\frac1{40},\qquad \mu_f=\kappa.}
\end{equation}
In particular, the all-character matrix
\eqref{JC:eq:regular-action-mode-matrix} is positive semidefinite.
\end{proposition}

\begin{proof}
The connection between the certified determinant and the Schur row is
division-free.  The limiting ledgers give
\[
 \bar g=2d_\circ^2Q_{\rm mat},\qquad
 \bar g\,\overline{\bar Z}
 =2d_\circ^2a e^{-i\pi x}\sin(\pi x)R_I,
\]
\[
 \bar A=-d_\circ^2a e^{-i\pi x}\sin(\pi x)R_{b,{\rm poly}},
 \qquad
 \bar r=d_\circ^2D_{\rm poly}(x;\mu_s).
\]
Consequently the limiting part of
\eqref{JC:eq:regular-action-reduced-Schur} after deletion of the fan reserve
is exactly
\begin{equation}\label{JC:eq:regular-proxy-Schur-determinant}
 \left(\frac12-\mu_s\right)\bar g|\bar Z|^2
 -\frac{|\bar A+\mu_s\bar g\,\overline{\bar Z}|^2}{\bar r}
 =\frac{d_\circ^2a^2\sin^2(\pi x)}
 {Q_{\rm mat}D_{\rm poly}(x;\mu_s)}\,\Delta_{\mu_s}(x).
\end{equation}
Both denominators are positive by
Proposition~\ref{JC:prop:regular-support-double-disk-gap}.

The support theorem gives
$p^{ss}_{N,k}\ge(1/2+1/20)g_{N,k}$, hence the denominator in
\eqref{JC:eq:regular-action-reduced-Schur}, with $\mu_s=1/40$, satisfies
$r_{N,k}\ge g_{N,k}/40>0$.  Suppose the claimed Schur row failed along a
sequence of characters.  After conjugation take $2\le k_N\le N/2$.
If $k_N$ is bounded, Lemma~\ref{JC:lem:regular-mixed-uniform-bridge} makes
the negative quotient and the corrector term $o(a^2)$; the untouched fan
term is $(3\kappa-\mu_f)a^2=2\kappa a^2$, a contradiction.

Otherwise pass to a subsequence with $k_N\to\infty$.  The same lemma
compares the quotient, including its denominator and corrector, with the
limiting proxy up to $o(a^2)$.  Proposition~\ref{JC:prop:regular-completed-covector-ratio}
makes that proxy nonnegative for $\mu_s=1/40$, including sequences with
$x_N\downarrow0$; at $x_N\to1/2$ its possible zero is again paid by the
same $2\kappa a^2$ fan reserve.  Thus the actual row is
$2\kappa a^2-o(a^2)>0$, another contradiction.  Every character sequence
has one of these two alternatives, proving the assertion uniformly.
\end{proof}

\begin{remark}[Why stationary O--H does not prove JC--AR]
On an even regular fan take the alternating support vector
$z_i=(-1)^i$.  The exact polygonal reconstruction in the CT node gives
\[
 \|V_z\|_{H^{1/2}(\mathbb T;\mathbb R^2)}^2\asymp Na^{-2},
 \qquad
 \Gjp_{\alpha_*}(z)\asymp N.
\]
Thus the full-vector norm in stationary O--H is larger than the stable
normal-sine metric by the sharp factor $a^{-2}$.  Its error can therefore
be as large as
\[
 C\omega_{OH}(a)a^{-2}\Gjp_{\alpha_*}(z).
\]
The proved statement $\omega_{OH}(a)\to0$ supplies neither the rate
$O(a^2)$ nor a sufficiently small leading constant.  Moreover its radial
disk trace is not literally the normal-sine completed trace, and it
controls the full Hessian rather than the base-subtracted action root
\eqref{JC:eq:regular-action-root}.  Hence O--H cannot be used to assert either
\eqref{JC:eq:regular-action-support-row} or
\eqref{JC:eq:regular-action-mode-matrix}.
\end{remark}

\begin{contract}[JC--SP: remaining super-near completed-action modulus and output]
\label{JC:con:JC-SP}
Suppose the regular completed-action contract JC--AR holds.  Let
\begin{equation}\label{JC:eq:super-near-energy}
 \mathcal N_N(d,w):=
 a^2\sum_i d_i^2+\Gjp_{\alpha_*}(w),
 \qquad
 \mathcal N_N(d,w)\le Ma^5.
\end{equation}
For the exact completed chart of JC--AR also put
\begin{equation}\label{JC:eq:finite-completed-energy}
 \mathfrak E_N(d,w):=
 \Jfour(\alpha_*+d)+
 \Gjp_{\alpha_*+d}(J_{\alpha_*+d}w).
\end{equation}
The finite-coordinate row needed here is closed by
Proposition~\ref{JC:prop:super-near-bridge-active}: uniformly for fixed $M$,
\begin{equation}\label{JC:eq:super-near-energy-bridge}
 c_M\mathcal N_N(d,w)\le\mathfrak E_N(d,w)
 \le C_M\mathcal N_N(d,w)
\end{equation}
whenever either side is at most a fixed multiple of $a^5$.  This row
includes the moving metric, the $J_\alpha$ transport, the quotient gauges,
and their uniform constants.  The exact finite corrector cancels here
because $z-z_\alpha^D=J_\alpha w$; its nonlinear variation instead belongs
to the completed chart and Hessian modulus below.

The exact fourth-defect identity gives the fan part of this bridge and the
radius
\begin{equation}\label{JC:eq:super-near-radius}
 \frac{\|d\|_{\ell^\infty}}a+
 a^{-2}\Gjp_{\alpha_*}(w)^{1/2}
 \le C_Ma^{1/2}.
\end{equation}
Indeed,
\[
 \Jfour(\alpha_*+d)
 =\sum_i d_i^2(6a^2+4ad_i+d_i^2)
 \ge3a^2\sum_i d_i^2
\]
for $d_i\ge-a$ and $\sum_i d_i=0$; the last coefficient is sharp at
$d_i=-a$.  The support part of
\eqref{JC:eq:super-near-energy-bridge} is the moving-low-mode transversality
content of Proposition~\ref{JC:prop:super-near-bridge-active}; it is not
inferred from a bare interchange of the two meshes.
Moreover, applying the closed corrector estimate to
$\alpha_*+td$, dividing by $t^2$, and letting $t\to0$ gives
\begin{equation}\label{JC:eq:corrector-derivative-bound}
 \Gjp_{\alpha_*}(Z_Nd)\le Ca^2\sum_i d_i^2.
\end{equation}
Thus the initial physical support velocity
$\zeta=w+Z_Nd$ also satisfies
$\Gjp_{\alpha_*}(\zeta)\le C_Ma^5$.

Use the affine path in the exact completed coordinates
\begin{equation}\label{JC:eq:completed-coordinate-path}
 q_t=(td,tw),\qquad
 \alpha_t=\alpha_*+td,\qquad
 z_t=J_{\alpha_t}(\bar z_{\alpha_t}^D+tw),
 \qquad0\le t\le1.
\end{equation}
Proposition~\ref{JC:prop:super-near-bridge-active} proves that this path
remains in the active convex chart and that all chart maps above are
defined uniformly.  Its Hessian is taken after
composition with the exact scalar \(\mathbf A_N\); hence the moving fan,
support, corrector, gauge, and coordinate-acceleration terms are all part
of the Hessian.  The earlier proposed affine chord of physical vertices
has zero vertex acceleration, but its initial chart tangent need not equal
the endpoint difference $(d,w)$.  It is therefore not used here without a
separate tangent-energy comparison.

For later estimates the exact completed scalar has the following
corrector-free normal form.  Put $c_\alpha=\bar z_\alpha^D$, so that
$z_\alpha^D=J_\alpha c_\alpha$.  Then
\begin{equation}\label{JC:eq:completed-action-corrector-normal-form}
 \boxed{
 \begin{aligned}
 \mathbf A_N(d,w)={}&
 \Phi_\alpha\!\left(h_\alpha\mathbf1
       +J_\alpha(c_\alpha+w)\right)
 -\Phi_\alpha(h_\alpha\mathbf1)\\
 &+\Gjp_\alpha(J_\alpha c_\alpha)
 -\frac12\Gjp_\alpha(J_\alpha w)
 +\frac\kappa2\Jfour(\alpha),
 \qquad \alpha=\alpha_*+d.
 \end{aligned}}
\end{equation}
Indeed $p_\alpha[\xi]=-2g_\alpha(z_\alpha^D,\xi)$, so the forcing,
metric, and completed square combine exactly into
$\Gjp_\alpha(z_\alpha^D)-\frac12\Gjp_\alpha(J_\alpha w)$.
Thus~\eqref{JC:eq:completed-action-corrector-normal-form} retains the moving
fan, support chart, quotient gauge, corrector, base-value subtraction,
moving metric, $1/2$ normalization, and fourth-defect row in one scalar.

The only remaining analytic row of JC--SP is a relative modulus for the
Hessian of the completed action itself.  For every fixed $M$ there must be
$\varepsilon_M(a)\to0$ such that
\begin{equation}\label{JC:eq:super-near-path-modulus}
 \boxed{
 \left|D^2\mathbf A_N(q_t)[(d,w),(d,w)]
       -D^2\mathbf A_N(0)[(d,w),(d,w)]\right|
 \le\varepsilon_M(a)\mathcal N_N(d,w)
 \quad(0\le t\le1).}
\end{equation}
A bound $\varepsilon_M(a)=O_M(a^{1/2})$ would be sufficient.  The finite
rank scale/translation gauge change, moving corrector, and base-value
subtraction are part of the left side and must not be dropped.

One stronger uniform operator formulation, sufficient for
\eqref{JC:eq:super-near-path-modulus}, is the following.  Let $R_N$ be the Riesz operator of
$\mathcal N_N$ on the fixed regular quotient and put
$H_N(q)=D^2\mathbf A_N(q)$.  With
\[
 r_N(q)=\frac{\|d(q)\|_{\ell^\infty}}a
 +a^{-2}\Gjp_{\alpha_*}(w(q))^{1/2},
\]
the natural operator statement is
\begin{equation}\label{JC:eq:super-near-operator-modulus}
 \boxed{
 \Omega_N(r):=
 \sup_{r_N(q)\le r}
 \left\|R_N^{-1/2}
 \{H_N(q)-H_N(0)\}R_N^{-1/2}\right\|_{\rm op}
 \longrightarrow0}
\end{equation}
as $r\downarrow0$, uniformly in $N$.  A Lipschitz estimate
$\Omega_N(r)\le Cr$, together with
$r_N(q_t)\le C_Ma^{1/2}$ from the chart/energy bridge, would imply the
stated $O_M(a^{1/2})$ rate.  This operator statement is stronger than the
diagonal path estimate and is recorded only as a sufficient formulation.
In block
form it must simultaneously contain
\begin{align*}
 |\Delta H_{ff}[d_1,d_2]|
 &\le Cr\,a^2\|d_1\|_2\|d_2\|_2,\\
 |\Delta H_{fs}[d,u]|
 &\le Cr\,a\|d\|_2\Gjp_{\alpha_*}(u)^{1/2},\\
 |\Delta H_{ss}[u_1,u_2]|
 &\le Cr\,
 \Gjp_{\alpha_*}(u_1)^{1/2}
 \Gjp_{\alpha_*}(u_2)^{1/2}.
\end{align*}
The nonzero period-three mixed row is therefore retained rather than
estimated as a separate remainder.

The absolute analytic CT estimate is not
\eqref{JC:eq:super-near-path-modulus}.  On the alternating support mode it
again loses $a^{-2}$, so a generic full-vector modulus gives only
$a^{-2}\omega(C_Ma^{1/2})\mathcal N_N$.  Analytic Lipschitz continuity or
$\omega(r)\to0$ does not make this quantity small.  The remaining part of
JC--SP therefore is an open natural-operator statement.

There is a scalar reduction which is weaker than proving the full operator
modulus.  For $q=(d,w)$ define
\[
 f_q(\zeta):=\mathbf A_N(\zeta d,\zeta w),\qquad
 r_N(q):=\frac{\|d\|_{\ell^\infty}}a
 +a^{-2}\Gjp_{\alpha_*}(w)^{1/2}.
\]
Suppose there are absolute $c_0,C_0>0$ such that every exact completed
coordinate in~\eqref{JC:eq:completed-action-corrector-normal-form} extends
holomorphically to $|\zeta|r_N(q)\le c_0$ and
\begin{equation}\label{JC:eq:super-near-complex-ray-bound}
 \boxed{|f_q''(\zeta)|\le C_0\mathcal N_N(d,w)}
\end{equation}
there.  If $r_N(q)\le c_0/2$, Cauchy's estimate on the disk of radius
$c_0/(2r_N(q))$ about $t\in[0,1]$ gives
\[
 |f_q'''(t)|\le\frac{2C_0}{c_0}r_N(q)\mathcal N_N(d,w).
\]
Since the completed-coordinate path is affine,
$f_q''(t)=D^2\mathbf A_N(q_t)[q,q]$.  Integration therefore proves
\begin{equation}\label{JC:eq:super-near-Cauchy-output}
 \boxed{
 |D^2\mathbf A_N(q_t)[q,q]-D^2\mathbf A_N(0)[q,q]|
 \le\frac{2C_0}{c_0}r_N(q)\mathcal N_N(d,w).}
\end{equation}
The bridge gives $r_N(q)=O_M(a^{1/2})$, so
\eqref{JC:eq:super-near-complex-ray-bound} would close the required path
modulus with $\varepsilon_M(a)=O_M(a^{1/2})$.

The precise analytic blocker behind
\eqref{JC:eq:super-near-complex-ray-bound} is a nonstationary quasiuniform
resummed completed-jet identity, abbreviated QRCI.  It must be obtained by
differentiating the exact scalar
\eqref{JC:eq:completed-action-corrector-normal-form}; the stage-130
polygon-minus-disk longitudinal ledger is retracted and is not such a
derivative.  Terminal two-ray reduced-DtN, its full-cell complement, the
negative cotangent/area pole, and every endpoint-localized derivative of
$h_\alpha,J_\alpha,c_\alpha$, the moving metric, and the base subtraction
must be combined before absolute values are taken.  The existing
full-vector Hessian estimate loses $a^{-2}$ on alternating support data,
while the available mesh-free completed estimate is stationary and uses a
degenerate radial comparison form.  Neither proves QRCI.

\begin{remark}[The scalar radial-chart shortcut fails at the critical index]
\label{JC:rem:super-near-radial-Hhalf-obstruction}
One cannot obtain the missing bound by differentiating only the scalar
radial graph.  If the fixed-normal support numbers move as
$h_i(t)=h_i+tz_i$, then on side $i$
$R_t(\varphi)=(h_i+tz_i)/\cos(\varphi-\theta_i)$.  At a regular vertex
the two one-sided traces of the formal derivative differ by
\begin{equation}\label{JC:eq:super-near-radial-Hhalf-jump}
 J_i=\frac{z_{i+1}-z_i}{\cos(a/2)}.
\end{equation}
Every nonconstant cyclic character has a nonzero such jump.  A jump is
not in $H^{1/2}(\mathbb T)$, because its Slobodetskii integral contains
\[
 \frac{|J_i|^2}{4}
 \int_0^{\varepsilon/2}\!\int_0^{\varepsilon/2}
 \frac{du\,dv}{(u+v)^2}=+\infty.
\]
Thus $t\mapsto R_t$ is not differentiable as an $H^{1/2}$-valued curve,
so the scalar CT analytic-Hessian theorem cannot be invoked.  Freezing
the vertices by a continuous physical affine-side field removes the jump
but restores the already recorded $a^{-2}$ loss.  This leaves the
endpoint-resummed NQCH cancellation as a genuine analytic input.
\end{remark}

\begin{remark}[Exact radius resummation, retraction of DAR, and the live completed third jet]
\label{JC:rem:super-near-completed-third-jet}
The artificial endpoint radius is not itself an obstruction.  If a
homogeneous full-cell corner power has amplitude $A$, exponent $\beta$, and
the terminal window has relative radius $\vartheta$, then in either the
common or jump channel, with $U_\beta$ and $D_\beta$ denoting the exact
two-ray common and jump coefficients,
\[
 A^2\vartheta^{2\beta}K_\beta
 +A^2(1-\vartheta^{2\beta})K_\beta=A^2K_\beta,
 \qquad K_\beta\in\{U_\beta,D_\beta\}.
\]
Hence the apparent $\beta\log\beta$ loss created by stopping at
$\vartheta=\beta$ cancels exactly, including after differentiation along
the completed ray.

The former amplitude row DAR is retracted as an independent contract.  It
was written, with
\[
 K_i(t)=\frac{D_{\beta_i(t)}}\pi j_i(q;t)^2,
\]
as $\sum_i(A_i^2)'K_i+(\sum_iC_i^{\rm end})'$.  But the exact endpoint
jet is unchanged by the analytic reallocation
\[
 \widetilde A_i^2=A_i^2+\gamma_i,
 \qquad
 \widetilde C_i^{\rm end}=C_i^{\rm end}-\gamma_iK_i,
\]
whereas the displayed DAR expression changes by
$-\sum_i\gamma_iK_i'$.  Choosing a sufficiently small constant nonzero
$\gamma_i$ and a nonconstant $K_i$ makes this change nonzero while
preserving $\widetilde A_i^2>0$ locally whenever $A_i^2>0$.  Thus the old
row depends on how the full-cell amplitude is
separated from its endpoint descendants; it is not a gauge-invariant
statement about the scalar~\eqref{JC:eq:completed-action-corrector-normal-form}.

The replacement live row is the third derivative of that scalar itself.
For $q=(d,w)$ and $q_t=tq$, put
$f_q(t)=\mathbf A_N(q_t)$.  The required real estimate is
\begin{equation}\label{JC:eq:super-near-completed-third-jet}
 \boxed{
 |f_q'''(t)|
 =|D^3\mathbf A_N(q_t)[q,q,q]|
 \le C r_N(q)\mathcal N_N(q),\qquad0\le t\le1.}
\end{equation}
It is gauge invariant and includes the full-cell flux, both terminal
channels, their exact complements, every endpoint and nonendpoint
chain-rule descendant, the moving metric, the corrector, the quotient
transport, and the base subtraction.

For clarity, it is obtained directly from the normal form.  If
\[
 Y_t=h_{\alpha_t}\mathbf1+J_{\alpha_t}(c_{\alpha_t}+tw),\quad
 B_t=h_{\alpha_t}\mathbf1,\quad
 C_t=J_{\alpha_t}c_{\alpha_t},\quad
 W_t=tJ_{\alpha_t}w,
\]
and $X_t=(\alpha_t,Y_t)$, then for every scalar family $F$
\[
 \frac{d^3}{dt^3}F(X_t)
 =D^3F(X_t)[X_t',X_t',X_t']
  +3D^2F(X_t)[X_t',X_t'']+DF(X_t)[X_t'''].
\]
Apply this identity to the two $\Phi$ terms and the two $\Gjp$ terms in
\eqref{JC:eq:completed-action-corrector-normal-form}, with $Y_t,B_t,C_t,W_t$
respectively, and add the exact fourth-defect derivative.  In particular
\[
 \frac\kappa2\frac{d^3}{dt^3}\Jfour(\alpha_*+td)
 =12\kappa\left(a\sum_i d_i^3+t\sum_i d_i^4\right),
\]
whose absolute value is at most
$Cr_N(q)a^2\sum_i d_i^2$.  No stage-130 ledger enters this chain rule.

The complex-ray bound~\eqref{JC:eq:super-near-complex-ray-bound} implies
\eqref{JC:eq:super-near-completed-third-jet} by the Cauchy estimate already
displayed above.  Conversely,
\eqref{JC:eq:super-near-completed-third-jet} directly implies the real path
modulus~\eqref{JC:eq:super-near-path-modulus} by integration, but does not by
itself assert a holomorphic-tube bound.  Either route would suffice for
JC--SP.  Neither is proved here; renaming the missing estimate does not
close QRCI or the main theorem.
\end{remark}

Its required output is nevertheless exact.  Put
$\mu_0=\min\{\mu_f,\mu_s\}$.  Taylor's formula in the completed chart gives
\[
 \mathbf A_N(d,w)=
 \int_0^1(1-t)D^2\mathbf A_N(q_t)[(d,w),(d,w)]\,dt.
\]
Thus JC--AR and~\eqref{JC:eq:super-near-path-modulus} imply, for sufficiently
small $a$ depending only on $M$,
\begin{equation}\label{JC:eq:JC-SP-output}
 \boxed{
 \mathbf A_N(d,w)\ge
 \left(\mu_0-\frac12\varepsilon_M(a)\right)\mathcal N_N(d,w)
 \ge0.}
\end{equation}
Together with the energy bridge~\eqref{JC:eq:super-near-energy-bridge}, this
is precisely the super-near action statement
\eqref{JC:eq:super-near-action-bound}; it is an output of JC--SP, not a
fourth assumed row.
\end{contract}

\begin{proposition}[Canonical finite bridge and completed-path activity]
\label{JC:prop:super-near-bridge-active}
There is a canonical analytic choice of the scale/translation
trivialization $J_\alpha$ near the regular fan for which
\eqref{JC:eq:super-near-energy-bridge} holds, with constants independent of
$N$, whenever either side is at most $Ma^5$.  For this choice the whole
path~\eqref{JC:eq:completed-coordinate-path} remains in the active convex
chart, and $J_{\alpha_t}$, its inverse, and the exact corrector are defined
uniformly for $0\le t\le1$.
\end{proposition}

\begin{proof}
Write $\alpha_i=a+d_i$, $\sum_i d_i=0$, and
\(\varrho=\|d\|_{\ell^\infty}/a\).  Choose normal angles with
$\theta_0=0$ and define the material map, for $0\le s\le a$, by
\[
 F_d(ia+s)=\theta_i+\frac{\alpha_i}{a}s .
\]
Thus
\begin{equation}\label{JC:eq:material-map-small}
 \|F_d'-1\|_{L^\infty}\le\varrho,
 \qquad
 \|F_d-\operatorname{id}\|_{L^\infty}\le2\pi\varrho .
\end{equation}
Composition by $F_d$ and $F_d^{-1}$ is consequently bounded on
$H^{1/2}(\mathbb T)$ uniformly when $\varrho\le1/2$, as follows directly
from the Gagliardo double-integral norm.

Direct integration of the two sine basis functions gives
\begin{equation}\label{JC:eq:exact-scale-functional}
 L_\alpha z:=\sum_i\left\{
 \tan\frac{\alpha_{i-1}}2+\tan\frac{\alpha_i}2\right\}z_i
 =\int_{\mathbb T}\mathcal I_\alpha z .
\end{equation}
The side with normal $\theta_i$ in the tangential reference polygon has
length $h_\alpha\{\tan(\alpha_{i-1}/2)+\tan(\alpha_i/2)\}$, so
$L_\alpha z=0$ is exactly the affine scale row.  For $b\in\mathbb R^2$
put $\tau_\alpha(b)_i=b\cdot e_r(\theta_i)$.  Since first harmonics solve
$f''+f=0$ on each cell,
\begin{equation}\label{JC:eq:exact-translation-interpolation}
 \mathcal I_\alpha\tau_\alpha(b)=b\cdot e_r
\end{equation}
exactly.

Let $\mathsf X_*$ be the regular nodal complement of characters
$0,1,N-1$.  For $w\in\mathsf X_*$ define
\[
 c_\alpha(w)=\frac{L_\alpha w}{L_\alpha\mathbf1},
 \qquad
 f_\alpha(w)=\mathcal I_\alpha(w-c_\alpha(w)\mathbf1),
\]
let $\beta_\alpha(w)$ be the $L^2$ projection coefficient of $f_\alpha(w)$
onto $\{\cos\theta,\sin\theta\}$, and set
\begin{equation}\label{JC:eq:canonical-Jalpha}
 J_\alpha w=w-c_\alpha(w)\mathbf1-\tau_\alpha(\beta_\alpha(w)).
\end{equation}
Equations~\eqref{JC:eq:exact-scale-functional} and
\eqref{JC:eq:exact-translation-interpolation} show that $J_\alpha w$ lies
in the moving scale/translation quotient and that $J_{\alpha_*}=I$.

We prove uniform transversality.  On a material cell, with $r=s/a$, the
two pulled-back basis functions are
\[
 B_0(\alpha_i,r)=\frac{\sin((1-r)\alpha_i)}{\sin\alpha_i},
 \qquad
 B_1(\alpha_i,r)=\frac{\sin(r\alpha_i)}{\sin\alpha_i}.
\]
Taylor expansion, cellwise integration, and real interpolation give, for
$\|v\|_{\ell_a^2}^2=a\sum_i|v_i|^2$,
\begin{align}
 \|\mathcal I_\alpha v\circ F_d-\mathcal I_*v\|_{L^2}
 &\le C\varrho a^2\|v\|_{\ell_a^2},\label{JC:eq:material-L2-error}\\
 \|\mathcal I_\alpha v\circ F_d-\mathcal I_*v\|_{H^{1/2}}
 &\le C\varrho a^{3/2}\|v\|_{\ell_a^2}.
 \label{JC:eq:material-Hhalf-error}
\end{align}
Indeed the pointwise basis error is $O(\varrho a^2)$ and its physical
derivative is $O(\varrho a)$.

Cyclic equivariance implies that the continuous Fourier support of the
regular sine interpolant of a nodal character $k$ lies in
$k+N\mathbb Z$.  Hence $\mathcal I_*w$ is exactly orthogonal to
$1,\cos\theta,\sin\theta$, not merely asymptotically so.  The regular
mass matrix and the disk multiplier give
\begin{equation}\label{JC:eq:regular-mass-G-control}
 \|w\|_{\ell_a^2}\le C\Gjp_{\alpha_*}(w)^{1/2}.
\end{equation}
Since
$|\tan(\alpha_i/2)-\tan(a/2)|\le C\varrho a$ and
$L_\alpha\mathbf1\asymp1$, weighted Cauchy--Schwarz yields
\begin{equation}\label{JC:eq:scale-coefficient-control}
 |c_\alpha(w)|\le C\varrho\Gjp_{\alpha_*}(w)^{1/2}.
\end{equation}
Changing variables $\theta=F_d(s)$ in the first-harmonic coefficients,
subtracting the exactly zero regular coefficient, and using
\eqref{JC:eq:material-map-small}--\eqref{JC:eq:scale-coefficient-control} gives
\begin{equation}\label{JC:eq:translation-coefficient-control}
 |\beta_\alpha(w)|\le C\varrho\Gjp_{\alpha_*}(w)^{1/2}.
\end{equation}
No inverse estimate enters this step: the multiplier
$e_r(F_d)F_d'-e_r$ is $O(\varrho)$ in $L^\infty$.

Put $T_\alpha w=\mathcal I_\alpha J_\alpha w$.  The functions
$(\mathcal I_\alpha\mathbf1)\circ F_d$ and $e_r\circ F_d$ have uniformly
bounded $H^{1/2}$ norms.  Therefore
\eqref{JC:eq:material-Hhalf-error},
\eqref{JC:eq:scale-coefficient-control}, and
\eqref{JC:eq:translation-coefficient-control} imply
\begin{equation}\label{JC:eq:moving-low-mode-transversality}
 \|T_\alpha w\circ F_d-\mathcal I_*w\|_{H^{1/2}}
 \le C\varrho\Gjp_{\alpha_*}(w)^{1/2}.
\end{equation}
Now $T_\alpha w$ is exactly orthogonal to the three disk null modes, so
$q_D(T_\alpha w)\asymp\|T_\alpha w\|_{H^{1/2}}^2$.  Absorbing the right
side of~\eqref{JC:eq:moving-low-mode-transversality} for a universal
$\varrho\le\varrho_0$ proves
\begin{equation}\label{JC:eq:moving-metric-equivalence}
 c\Gjp_{\alpha_*}(w)\le
 \Gjp_\alpha(J_\alpha w)\le C\Gjp_{\alpha_*}(w).
\end{equation}
It also proves that $J_\alpha$ is injective.  The two quotient spaces
have dimension $N-3$, hence it is onto and has a uniform inverse.

The exact identity
\[
 \Jfour(\alpha_*+d)
 =\sum_i d_i^2(6a^2+4ad_i+d_i^2)
\]
has coefficient at least $3a^2$ for every $d_i\ge-a$.  If either side of
\eqref{JC:eq:super-near-energy-bridge} is at most $Ma^5$, it follows first
that $\varrho\le C_Ma^{1/2}$, hence $\varrho\le\varrho_0$ for large $N$.
The same coefficient is then bounded above by a universal multiple of
$a^2$.  Combining this fact with
\eqref{JC:eq:moving-metric-equivalence} proves
\eqref{JC:eq:super-near-energy-bridge}.  The corrector does not occur in
this comparison because $z-z_\alpha^D=J_\alpha w$ identically.

It remains to prove activity.  If $\Lambda_{\alpha,i}(z)$ denotes the
change of side $i$ under a support increment $z$, the exact support-line
formula is
\begin{equation}\label{JC:eq:exact-side-length-jump}
 \Lambda_{\alpha,i}(z)=
 \frac{z_{i+1}-z_i\cos\alpha_i}{\sin\alpha_i}
 +\frac{z_{i-1}-z_i\cos\alpha_{i-1}}{\sin\alpha_{i-1}}
 =f'(\theta_i+)-f'(\theta_i-),
 \quad f=\mathcal I_\alpha z.
\end{equation}
After subtracting an exact translation, which changes no side length,
the quasiuniform one-dimensional finite-element inverse estimates give
\[
 \|f\|_{H^1}\le Ca^{-1/2}\|f\|_{H^{1/2}},
 \qquad
 \|f'\|_{L^\infty}\le Ca^{-1/2}\|f'\|_{L^2}.
\]
Because $L_\alpha z=0$ removes the constant mode, the disk multiplier is
equivalent to the remaining $H^{1/2}$ norm.  Thus
\begin{equation}\label{JC:eq:side-length-inverse}
 \max_i|\Lambda_{\alpha,i}(z)|
 \le Ca^{-1}\Gjp_\alpha(z)^{1/2}.
\end{equation}

Along~\eqref{JC:eq:completed-coordinate-path}, exact linearity gives
$z_t=z_{\alpha_t}^D+J_{\alpha_t}(tw)$.  The closed corrector estimate,
the fourth-defect identity, and
\eqref{JC:eq:moving-metric-equivalence} give
\[
 \Gjp_{\alpha_t}(z_{\alpha_t}^D)\le C_Ma^5,
 \qquad
 \Gjp_{\alpha_t}(J_{\alpha_t}(tw))\le C_Ma^5.
\]
Consequently $\Gjp_{\alpha_t}(z_t)^{1/2}\le C_Ma^{5/2}$.
The reference side lengths are at least $ca$, whereas
\eqref{JC:eq:side-length-inverse} bounds every change by
$C_Ma^{3/2}=o(a)$.  Thus every side remains positive.  The companion
nodal and derivative inverse bounds give normal and tangential vertex
displacements $O_M(a^2)$ and $O_M(a^{3/2})$, respectively, so the whole
path remains in a fixed smaller radial $W^{1,\infty}$ chart.  This proves
all assertions.
\end{proof}

\begin{remark}[Scope of the regular/super-near split]
JC--AR and JC--SP isolate the regular completed-action root, the finite
energy bridge, and its diagonal propagation.  Their exact output is only
the super-near action bound~\eqref{JC:eq:JC-SP-output}.  The full no-ratio
Contract~\ref{JC:con:JC-PL} additionally requires JC--CL to absorb every
endpoint outside~\eqref{JC:eq:super-near-energy}.  No far implication is
assumed in JC--AR or JC--SP.
\end{remark}

\begin{contract}[JC--CL: coarse/far action localization]
\label{JC:con:JC-CL}
Write $d=d_\alpha=z_\alpha^D$ and $z=d+w$.  Exact disk completion first
removes every disk covector from the coarse problem:
\begin{equation}\label{JC:eq:coarse-action-normal-form}
 \boxed{
 \mathscr A_\alpha(d+w)=
 \Phi_\alpha(h_\alpha\mathbf1+d+w)
 -\Phi_\alpha(h_\alpha\mathbf1)
 -\frac12\Gjp_\alpha(w)+\Gjp_\alpha(d)
 +\frac\kappa2\Jfour(\alpha).}
\end{equation}
Indeed
$q_D(B_\alpha+K_\alpha z)-q_D(B_\alpha)
=\Gjp_\alpha(z-d)-\Gjp_\alpha(d)$, and
Proposition~\ref{JC:prop:JC-ALG} gives the formula.  This identity also fixes
the admissible path.  The active chart supplies only
$h_\alpha\mathbf1+[0,1](d+w)$; it does not imply that
$h_\alpha\mathbf1+d$ or
$h_\alpha\mathbf1+d+[0,1]w$ is active.  In the presence of an arbitrarily
short side, the energy bound $\Gjp_\alpha(d)\lesssim\Jfour(\alpha)$ has no
ratio-free inverse side-length consequence.  Thus a coarse proof may not
Taylor-expand about the disk corrector without a new activity lemma.

Put
\begin{equation}\label{JC:eq:coarse-action-energy}
 E_\alpha(z):=\Jfour(\alpha)
 +\Gjp_\alpha(z-d_\alpha).
\end{equation}
A sufficient coarse estimate is the existence of constants
$c_F,C_F>0$ such that, on the whole active no-ratio chart,
\begin{equation}\label{JC:eq:coarse-action-bound}
 \boxed{
 \mathscr A_\alpha(z)
 \ge c_FE_\alpha(z)-C_Fa^5.}
\end{equation}
For $E_\alpha(z)\le Ma^5$, set
$d=\alpha-\alpha_*$ and
$w=J_\alpha^{-1}(z-d_\alpha)$.  The energy bridge and output required in
JC--SP then give
\begin{equation}\label{JC:eq:super-near-action-bound}
 \mathscr A_\alpha(z)\ge0
 \quad\hbox{whenever}\quad
 E_\alpha(z)\le Ma^5
\end{equation}
for every fixed $M$.  Taking $M\ge C_F/c_F$, Contract~\ref{JC:con:JC-PL}
follows.  Indeed,
\eqref{JC:eq:super-near-action-bound} handles the first branch, while on
$E_\alpha(z)>Ma^5$ equation~\eqref{JC:eq:coarse-action-bound} gives
\[
 \mathscr A_\alpha(z)
 >(c_FM-C_F)a^5\ge0.
\]

No retained estimate proves~\eqref{JC:eq:coarse-action-bound}.  The existing
two-regime estimates for $F_4$, the normal-sine/Bessel/secant transfers,
and JE--V concern pure fan disk endpoints or scalar tails.  They do not
control arbitrary support amplitudes.  In particular, at
$\alpha=\alpha_*$ one may have
$\Jfour=0$ and $\Gjp_{\alpha_*}(z)\gg a^5$ while remaining in the radial
chart.  The fixed-fan Hessian identity JP--EH supplies no lower bound for
that branch.  The nonzero physical-curvature mixed row also cannot be
inserted as a vanishing error, by the CT mixed counterexample.

This fixed-fan obstruction has an exact scalar form.  Cyclic symmetry at
$\alpha_*$ gives $b_*=d_*=0$, so JC--ALG reduces to
\[
 \mathscr A_*(z)=
 \Phi_*(h_*\mathbf1+z)-\Phi_*(h_*\mathbf1)
 -\frac12\Gjp_*(z).
\]
Consequently the restriction of~\eqref{JC:eq:coarse-action-bound} to the
regular fan requires the following regular-fan Bregman coercivity row,
with $c_s=c_F$ and $C=C_F$:
\begin{equation}\label{JC:eq:regular-fan-Bregman-coercivity}
 \boxed{
 \Phi_*(h_*\mathbf1+z)-\Phi_*(h_*\mathbf1)
 \ge\left(\frac12+c_s\right)\Gjp_*(z)-Ca^5.}
\end{equation}
Equivalently, JP--EH rewrites its left side as the support-time integral
of the exact regular-fan reduced Helmholtz DtN, area, and lower rows.  The
two classical Minkowski concavities for $|P|^{1/2}$ and
$\lambda_1(P)^{-1/2}$ do not order their ratio and hence do not prove this
inequality.  Thus~\eqref{JC:eq:regular-fan-Bregman-coercivity}, followed by
the genuinely nonregular coarse localization, is the smallest presently
isolated content of JC--CL.

At the regular fan cyclic symmetry gives $\ell_*=d_*=0$.  Define
\begin{equation}\label{JC:eq:regular-fan-integrated-modulus}
 \mathcal L_*(z):=\int_0^1(1-t)
 \{H_{*,tz}-H_{*,0}\}[z,z],dt.
\end{equation}

\begin{proposition}[The fixed-reserve RF--IM target is false]
\label{JC:prop:RF-IM-period-three-counterexample}
Let \(N=3M\to\infty\), \(a=2\pi/N\), and, on the regular fan, put
\[
 d_i=\cos\frac{2\pi i}{3},\qquad
 s_*=\frac{4h_*\sin^2(a/2)}{1+2\cos a},\qquad
 z_N=(1-a^2)s_*d.
\]
The segment \(h_*\mathbf1+[0,1]z_N\) is strictly active and remains in
the radial small chart.  Moreover
\begin{align}
 \Gjp_*(z_N)
 &=\frac{13d_\circ^2\zeta(3)}{4\pi^2}a^3\{1+o(1)\},
 \label{JC:eq:RF-counterexample-G}\\
 \frac{\Phi_*(h_*\mathbf1+z_N)-\Phi_*(h_*\mathbf1)}
 {\Gjp_*(z_N)}
 &\longrightarrow\frac{10}{13},
 \label{JC:eq:RF-counterexample-value}\\
 \frac{\frac12H_{*,0}[z_N,z_N]}{\Gjp_*(z_N)}
 &\longrightarrow\frac{2\pi^2\log3}{13\zeta(3)}.
 \label{JC:eq:RF-counterexample-root}
\end{align}
Consequently
\begin{equation}\label{JC:eq:regular-fan-integrated-modulus-disproved}
 \frac{\mathcal L_*(z_N)}{\Gjp_*(z_N)}
 \longrightarrow
 \frac{10}{13}-\frac{2\pi^2\log3}{13\zeta(3)}
 =-0.6185019068\ldots<-\frac1{20}.
\end{equation}
Thus no \(N\)-independent \(C\) and \(\eta<1/20\) can make the former
RF--IM bound
\(\mathcal L_*(z)\ge-\eta\Gjp_*(z)-Ca^5\) valid on every regular active
radial path.  This does not disprove RF--BC: on the same family
\begin{equation}\label{JC:eq:RF-counterexample-action}
 \frac{\mathscr A_*(z_N)}{\Gjp_*(z_N)}
 \longrightarrow\frac7{26}>0.
\end{equation}
\end{proposition}

\begin{proof}
For a fixed normal fan,
\[
 L_i(h)=\frac{h_{i-1}+h_{i+1}-2h_i\cos a}{\sin a}.
\]
Since \(d_{3j}=1\) and \(d_{3j+1}=d_{3j+2}=-1/2\), exactly the
\(3j\)-sides first vanish at \(s=s_*\), while all other sides grow.
Taking \(s_N=(1-a^2)s_*\) makes the entire segment strictly active.
Also \(s_*=a^2/3+O(a^4)\), so the support displacement is \(O(a^2)\)
and its adjacent difference quotient is \(O(a)\).

At \(s=s_*\), deletion of the zero sides leaves equal support numbers and
alternating normal gaps \(a,2a\).  Denote the area-normalized disk sources
of the regular and alternating fans by \(B_{\rm reg}\) and \(B_{\rm alt}\).
In the Fourier normalization used here, their leading material profiles
satisfy
\[
 F_{\rm reg}''=1-\sum_{j\in\mathbb Z}\delta_{j+1/2},
\qquad
 F_{\rm alt}''=1-\delta_{1/2}-2\delta_2
\]
with periods \(1\) and \(3\), respectively.  Hence
\[
 \|B_{\rm reg}\|_{\dot H^{1/2}}^2
 =\frac{\zeta(3)}{4\pi^3}a^3+O(a^4).
\]
For \(\varkappa=2\pi/3\),
\[
 \widehat F_{{\rm alt},\ell}
 =\frac{e^{-i\varkappa\ell/2}+2e^{-2i\varkappa\ell}}
 {3(\varkappa\ell)^2}.
\]
The numerator has squared modulus \(9\) for even \(\ell\) and \(1\) for
odd \(\ell\).  The even and odd \(\ell^{-3}\) sums are
\(\zeta(3)/8\) and \(7\zeta(3)/8\), so
\[
 \|B_{\rm alt}\|_{\dot H^{1/2}}^2
 =\frac{3\zeta(3)}{2\pi^3}a^3+O(a^4).
\]
The source-profile Taylor remainder is \(O(a^{7/2})\) in \(H^{1/2}\).
Since the first nonconstant frequencies are of order \(a^{-1}\) and the
disk multiplier is \(\omega_n=|n|+1+O(|n|^{-1})\),
\begin{equation}\label{JC:eq:RF-counterexample-source}
 q_D(B_{\rm alt})-q_D(B_{\rm reg})
 =\frac{5d_\circ^2\zeta(3)}{2\pi^2}a^3+O(a^4).
\end{equation}

The closed disk-chart scalar expansion applies separately to both
quasiuniform fans:
\[
 \Phi_\alpha(h_\alpha\mathbf1)
 =\lambda_\circ+q_D(B_\alpha)+O(a^4).
\]
Indeed the complete cubic term is bounded by
\(\|B\|_{W^{1,\infty}}\|B\|_{H^{1/2}}^2=O(a^4)\), and the tame
fourth-and-higher remainder is \(O(a^5)\).  The support distance from
\(s_N\) to the collapsed endpoint is \(O(a^4)\); two-sided parallel-body
inclusions and Dirichlet monotonicity change \(\Phi\) by \(O(a^4)\).
Thus~\eqref{JC:eq:RF-counterexample-source} gives the numerator in
\eqref{JC:eq:RF-counterexample-value}.

For the regular normal-sine metric,
\[
 Q(x)=\frac{2\sin^4(\pi x)}{\pi^3}
 \{\zeta(3,x)+\zeta(3,1-x)\}.
\]
The real character \(d\) has half the energy of the unnormalized complex
character.  Since
\[
 \zeta(3,1/3)+\zeta(3,2/3)=26\zeta(3),\qquad
 Q(1/3)=\frac{117\zeta(3)}{4\pi^3},
\]
and \(s_N^2=a^4/9+O(a^6)\), the factor
\((N/2)d_\circ^2Q(1/3)\) proves~\eqref{JC:eq:RF-counterexample-G}, and
then~\eqref{JC:eq:RF-counterexample-value}.

Finally the true polygonal support symbol gives
\[
 \frac{\frac12H_{*,0}[z_N,z_N]}{\Gjp_*(z_N)}
 \longrightarrow\frac{P(1/3)}{Q(1/3)}.
\]
The digamma third-point formula yields
\[
 P(1/3)=\frac{9\log3}{2\pi},\qquad
 \frac{P(1/3)}{Q(1/3)}
 =\frac{2\pi^2\log3}{13\zeta(3)},
\]
which proves~\eqref{JC:eq:RF-counterexample-root}.  The exact Taylor identity
\[
 \mathcal L_*(z_N)
 =\Phi_*(h_*\mathbf1+z_N)-\Phi_*(h_*\mathbf1)
 -\frac12H_{*,0}[z_N,z_N]
\]
now proves~\eqref{JC:eq:regular-fan-integrated-modulus-disproved}.  The strict
comparison with \(-1/20\) follows already from
\(\pi^2>9\), \(\log3>1\), and \(\zeta(3)<5/4\).  Since
\(\Gjp_*(z_N)\asymp a^3\), an \(Ca^5\) remainder cannot absorb it.
Subtracting \(1/2\) from the value ratio gives~\eqref{JC:eq:RF-counterexample-action}.
\end{proof}

\begin{contract}[RF--REL: relative preservation of the full root reserve]
\label{JC:con:RF-relative-root-reserve}
There exist \(\delta>0\) and \(C\), independent of \(N\), such that every
regular active radial path satisfies
\begin{equation}\label{JC:eq:RF-relative-root-reserve}
 \boxed{\mathcal L_*(z)\ge-\frac{1-\delta}{2}
 \{H_{*,0}[z,z]-\Gjp_*(z)\}-Ca^5.}
\end{equation}
\end{contract}

\begin{lemma}[RF--REL implies RF--BC]
\label{JC:lem:RF-relative-implies-Bregman}
Contract~\ref{JC:con:RF-relative-root-reserve} implies
\[
 \mathscr A_*(z)\ge\frac{\delta}{20}\Gjp_*(z)-Ca^5.
\]
\end{lemma}

\begin{proof}
Taylor's identity and the regular action give
\[
 \mathscr A_*(z)
 =\frac12\{H_{*,0}[z,z]-\Gjp_*(z)\}+\mathcal L_*(z).
\]
RF--REL and
\(H_{*,0}\ge(11/10)\Gjp_*\) imply
\[
 \mathscr A_*(z)
 \ge\frac{\delta}{2}\{H_{*,0}[z,z]-\Gjp_*(z)\}-Ca^5
 \ge\frac{\delta}{20}\Gjp_*(z)-Ca^5.
\]
\end{proof}

Thus the regular coarse row remains open, but its current sufficient
interface is RF--REL rather than the disproved fixed \(1/20\)-reserve
target.

\begin{remark}[Collapse-word audit and the exact JC--CL cut]
\label{JC:rem:RF-collapse-word-audit}
Every fixed periodic pattern of deleted regular-fan sides can be encoded
by a positive integer gap word $(m_1,\ldots,m_r)$.  At its collapse
endpoint the support profile on a gap of length $m$ is
$u_j=j(m-j)/2$, and the disk-source, metric, and root-reserve
coefficients reduce to finite Fourier/Hurwitz-zeta sums.  These formulas
recover exactly the period-three ratios $10/13$ and $7/26$ above.
The reproducible integer-word precheck
\texttt{diagnostics/explore\_regular\_collapse\_words.py} enumerates all
words of total period at most $16$ and several families through period
$4096$.  It finds no negative action and no vanishing relative reserve;
the smallest sampled ratios are approximately
\[
 \mathscr A_*/\Gjp_*=0.0912513,
 \qquad \mathscr A_*/\mathcal R_0=0.181689,
 \qquad \mathcal R_0=(H_{*,0}-\Gjp_*)/2.
\]
These numbers are not lower bounds on the full active microscopic class.
For $z_i=a^2u_i$, the exact leading envelope is
\[
 B_u(x)=\min_i\left\{u_i+\frac{(x-i)^2}{2}\right\},
 \quad s_i=i+\frac12+u_{i+1}-u_i,
 \quad g_i=s_i-s_{i-1}\ge0.
\]
At fixed microscopic period, the complete \emph{leading} active class is
the gap simplex
$g_i\ge0$, $\sum_i g_i=p$, and along support time
$g_i(t)=(1-t)+tg_i$.  If $S_p(t)$ is the resulting leading disk-source
energy, its exact pair ledger is
\[
 S_p''(t)=\frac1{\pi p}\sum_{i,j}(v_i-v_j)^2
 \log\left|2\sin\frac{\pi\{(i-j)+t(v_i-v_j)\}}p\right|,
 \qquad v_i=u_{i+1}-u_i.
\]
Here the diagonal terms are defined by \(0^2\log0:=0\).  If distinct
corner atoms collide at \(t=1\), this formula and the support-time Taylor
integral are understood as improper identities; the resulting logarithmic
singularity is integrable against \(1-t\).
Put
\[
 \widehat u_k=\frac1p\sum_{j=0}^{p-1}u_j e^{-2\pi i k j/p},
 \qquad
 \mathcal Q_p[u]=2\pi\sum_{k=1}^{p-1}
       |\widehat u_k|^2Q(k/p).
\]
Consequently the smallest leading regular target is the universal
monotone-transport inequality
\begin{equation}\label{JC:eq:RF-obstacle-BC}
 2\pi\{S_p(1)-S_p(0)\}
 \ge\left(\frac12+c_0\right)\mathcal Q_p[u]
\end{equation}
for every period and every point of that simplex.  Continuous-simplex
searches lower the sampled ratios to about $0.082293$ relative to
$\mathcal Q_p$ and $0.163711$ relative to the root reserve at $p=384$;
they find no negative value, but do not prove a uniform constant.
Equation~\eqref{JC:eq:RF-obstacle-BC}, its finite-$N$ secant/cubic/tail
transfer, and the nonregular original-path row all remain open.
This is a floating-point necessary-family precheck, not a proof of
RF--REL.  It also shows that the regular path loss is order one rather
than an $o(1)\Gjp_*$ error.  Even a proof of RF--REL would leave the
genuinely nonregular original-path inequality below as a second,
independent input.
\end{remark}

For a genuinely nonregular fan the exact original-path form is, with
$z=d+w$,
\begin{equation}\label{JC:eq:nonregular-original-path-action}
 \boxed{
 \begin{aligned}
 \mathscr A_\alpha(d+w)={}&\ell_\alpha[d+w]
 +\int_0^1(1-t)
 H_{\alpha,t(d+w)}[d+w,d+w],dt\\
 &-\frac12\Gjp_\alpha(w)+\Gjp_\alpha(d)
 +\frac\kappa2\Jfour(\alpha).
 \end{aligned}}
\end{equation}
The nonregular part of JC--CL is precisely a lower bound for this one
budget by $c_F\{\Jfour(\alpha)+\Gjp_\alpha(w)\}-C_Fa^5$, uniformly in
all gap ratios.  Separating its covector from its integrated Hessian would
recreate the unavailable JE--FM step.

For a stationary-minimizer bypass one may replace
\eqref{JC:eq:coarse-action-bound} by the weaker localization statement:
every stationary competitor $x$ with
$\Phi_N(x)\le\Phi_N(0)$ satisfies
\begin{equation}\label{JC:eq:stationary-coarse-localization}
 6a^2\sum_i d_i^2
 +\Gjp_{\alpha_*}(z-Z_Nd)\le Ma^5.
\end{equation}
Together with JC--AR and JC--SP, this would put the competitor in the
super-near action-positive branch and exclude a lower value through the
exact JP--JC/NR assembly.  The present global wrapper
gives only qualitative near-disk entry, not
\eqref{JC:eq:stationary-coarse-localization}.  Thus the stationary route has
three open analytic rows, not two.
\end{contract}

\begin{contract}[JC--PL: pathwise LCA payment]
\label{JC:con:JC-PL}
After decreasing the no-ratio fan and radial endpoint-chart thresholds so that
$\varepsilon_V(S)\le\kappa/2$, prove, for every active normalized support
vector $z\in\mathsf X_\alpha$ in that chart,
\begin{equation}\label{JC:eq:JC-PL}
 \boxed{\mathscr A_\alpha(z)\ge0.}
\end{equation}
The proof must use a disjoint, complete pair ledger.  The terminal two-ray
power and the negative cotangent/area jump row must be joined before any
estimate; every unused terminal annulus must remain in the nonterminal
ledger; same-cell, adjacent, and first-separation LCA pairs must be counted
exactly once.  The support-time integral must be taken before a side-ratio
supremum, using Lemma~\ref{JC:lem:support-time-log}.  The true ground-state
multiplier, reduced-energy, mixed-flux, area, and fixed-rank rows must
retain the same completed covector $b_\alpha$ of
\eqref{JC:eq:completed-covector}.  All constants must be independent of the
smallest angle and adjacent ratios.  In particular the contract includes
endpoints with an arbitrarily short positive side: smallness of an
affine-side parametrization relative to the tangential fan polygon is not
an assumption.  Its admissible chart is exactly the closed JP--NS output
\[
 h_\alpha\mathbf1+[0,1]z\ \hbox{active},\qquad
 \sup_{0\le t\le1}\|\rho_t\|_{W^{1,\infty}}
 +\Gjp_\alpha(z)^{1/2}\le C(S+\rho_0),
\]
not the false labelled affine-side chart of
Proposition JP--NS~\textup{(radial-not-affine)}.

Equivalently, by Proposition~\ref{JC:prop:JC-ALG}, prove
\begin{equation}\label{JC:eq:JC}
\boxed{
 \mathscr R_\alpha(z)
 \ge-\frac\kappa2\Jfour(\alpha)
      -\frac12\Gjp_\alpha(z-z_\alpha^D).}
\end{equation}
The stage-130 moving-disk longitudinal calculation omits material
chain-rule terms and is not an admissible proof of either display.
\end{contract}

\begin{theorem}[JC--PL closes the active-chart theorem]\label{JC:thm:JC-to-J0}
Assume Contract~\ref{JC:con:JC-PL}.  Then every normalized polygon in the active
near-disk chart satisfies
\begin{equation}\label{JC:eq:J0}
\boxed{
 \Phi_\alpha(h_\alpha\mathbf1+z)
 -\Phi_{\alpha_*}(h_*\mathbf1)
 \ge\kappa\Jfour(\alpha)
      +\frac12\Gjp_\alpha(z-z_\alpha^D).}
\end{equation}
Moreover
\begin{equation}\label{JC:eq:corrector-bound}
 \Gjp_\alpha(z_\alpha^D)
 \le\frac{C_p^2}{4}\Jfour(\alpha).
\end{equation}
\end{theorem}

\begin{proof}
Add and subtract the two disk endpoint values.  Equations
\eqref{JC:eq:joint-residual}, \eqref{JC:eq:disk-completion}, and \eqref{JC:eq:V}
give
\begin{align*}
 \Phi_\alpha(h_\alpha\mathbf1+z)
 -\Phi_{\alpha_*}(h_*\mathbf1)
 \ge{}&P_D^{\rm ns}(\alpha)-P_D^{\rm ns}(\alpha_*)
       +\Gjp_\alpha(z-z_\alpha^D)\\
 &-\varepsilon_V(S)\Jfour(\alpha)+\mathscr R_\alpha(z).
\end{align*}
Insert \eqref{JC:eq:D0}, \eqref{JC:eq:JC}, and
$\varepsilon_V(S)\le\kappa/2$.  The fan coefficient is
$2\kappa-\kappa/2-\kappa/2=\kappa$, and the support coefficient is
$1-1/2=1/2$.  This proves \eqref{JC:eq:J0}.  Equations \eqref{JC:eq:P} and
\eqref{JC:eq:corrector-size} prove \eqref{JC:eq:corrector-bound}.
\end{proof}

\begin{remark}[Precise mathematical cut]
Contract~\ref{JC:con:JC-PL} is not a restatement of either discarded separated
estimate.  At $\alpha=\alpha_*$ it asks only for the coarse one-sided
support bound encoded by \eqref{JC:eq:JC}; away from the regular fan it gives
the forcing access to the complete integrated Hessian surplus before a
Schur complement is taken.  The exact algebra and the support-time
logarithm cancellation are now proved.  Revision 24 further resolves the
former vague ``regular plus path'' description into one closed and two
open exact subcontracts.  The all-character regular completed-action Schur matrix
JC--AR is now closed.  The two remaining open contracts are the
completed-action Hessian modulus JC--SP and the coarse/far localization
JC--CL.  The former period-three completed-covector claim remains
retracted, and the fixed-reserve RF--IM target is disproved.  What is
still missing is the ratio-free
pathwise LCA payment for the true nonstationary ground-state multiplier and
all remaining lower rows.  No retained node implies either open analytic
estimate.
\end{remark}

\subsection{DAG disposition}

\begin{center}
\begin{longtable}{>{\raggedright\arraybackslash}p{.34\linewidth}
                  >{\raggedright\arraybackslash}p{.54\linewidth}}
\toprule
item & revision-24 disposition\\
\midrule
$D0_{\rm ns}$, JE--P, JE--V & closed inputs, equations
\eqref{JC:eq:D0}--\eqref{JC:eq:V};\\
JE--FM & unavailable; the old material falsification is retracted and the
row is absent from the active DAG;\\
JP--SC & retained only as a possible stronger auxiliary estimate; no
longer an active prerequisite;\\
JC--ALG and JC--TL & closed by Proposition~\ref{JC:prop:JC-ALG} and
Lemmas~\ref{JC:lem:support-time-log}--\ref{JC:lem:support-time-min};\\
pure-support JC--AR row & closed for all characters by
Lemma~\ref{JC:lem:fixed-character-support-limit} and
Proposition~\ref{JC:prop:regular-support-renormalized};\\
period-three completed-covector row & retracted by
Remark~\ref{JC:rem:period-three-completion-retracted}; its true polygonal
mixed finite part is now closed at every fixed Bloch ratio, as is the
strict support--double-disk gap; the ratio certificate and uniform
all-character bridge are closed by
Proposition~\ref{JC:prop:regular-action-Schur-closed};\\
RF--IM & disproved by
Proposition~\ref{JC:prop:RF-IM-period-three-counterexample}; RF--REL is the
replacement open sufficient interface for RF--BC;\\
super-near bridge/activity & closed by
Proposition~\ref{JC:prop:super-near-bridge-active};\\
terminal-radius pure-power resummation & closed by
Remark~\ref{JC:rem:super-near-completed-third-jet}; the old DAR row is
retracted as noninvariant, while the completed third-jet/QRCI row stays open;\\
JC--AR & closed all-character regular action Schur contract;\\
JC--SP, JC--CL & remaining natural Hessian-modulus and coarse/far
subcontracts inside JC--PL;\\
JC--PL & unique open conceptual node, Contract~\ref{JC:con:JC-PL};\\
JC--PL $\to$ JP--JC $\to$ J0 & closed algebraically by
Theorem~\ref{JC:thm:JC-to-J0}.\\
\bottomrule
\end{longtable}
\end{center}
\endgroup
